% !TEX program = xelatex
% Canonical source of the current manuscript. Historical paper files are retained.
\documentclass[UTF8,zihao=-4,fontset=fandol]{ctexart}
\usepackage[a4paper,margin=24mm,headheight=16pt]{geometry}
\usepackage{amsmath,amssymb,amsthm,mathtools,booktabs,longtable,array,tabularx}
\usepackage{enumitem,xcolor,fvextra,graphicx}
\usepackage[unicode,colorlinks=true,linkcolor=blue!45!black,citecolor=blue!45!black,urlcolor=blue!45!black]{hyperref}
\usepackage{xurl}
\usepackage{fancyhdr}
\pagestyle{fancy}\fancyhf{}
\fancyhead[L]{\small 相邻六次幂之间的素数}
\fancyhead[R]{\small 本地审阅稿}
\fancyfoot[C]{\thepage}
\setlength{\parindent}{2em}
\setlength{\parskip}{0.3em}
\setlength{\emergencystretch}{4em}
\setlist{itemsep=0.2em,topsep=0.4em}
\allowdisplaybreaks[2]
\newtheorem{theorem}{定理}[section]
\newtheorem{lemma}[theorem]{引理}
\newtheorem{proposition}[theorem]{命题}
\newtheorem{corollary}[theorem]{推论}
\theoremstyle{definition}
\newtheorem{definition}[theorem]{定义}
\theoremstyle{remark}
\newtheorem{remark}[theorem]{说明}
\renewcommand{\proofname}{证明}
\newcommand{\R}{\mathbb R}
\newcommand{\C}{\mathbb C}
\newcommand{\N}{\mathbb N}
\newcommand{\ii}{\mathrm i}
\newcommand{\e}{\mathrm e}
\newcommand{\dd}{\,\mathrm d}
\DeclareMathOperator{\Rea}{Re}
\DeclareMathOperator{\Ima}{Im}
\newcommand{\leanref}[2]{{\footnotesize\textsf{Lean：}\nolinkurl{#1.lean}\textsf{ :: }\nolinkurl{#2}}}
\newcommand{\repofile}[1]{\nolinkurl{#1}}
\hypersetup{pdftitle={相邻六次幂之间的素数：显式条件证明及 Lean 形式化},pdfauthor={Claude; Codex},pdfsubject={QRH(7/8), seven explicit inputs, conditional sixth-power prime intervals}}
\title{相邻六次幂之间的素数：\\显式条件证明及 Lean 形式化}
\author{Codex：QRH 路线研究、实现与本稿整理\\Claude：传统路线负责人；QRH 早期版本审阅}
\date{2026 年 11 月 6 日\quad 本地审阅稿}
\begin{document}
\maketitle
\begin{abstract}
在准黎曼条件 $\Rea\rho\le7/8$、固定版本的显式零点密度估计及五项逐一陈述的经典或计算文献输入下，证明对每个整数 $n\ge1$，存在素数 $p$ 满足
\[
n^6<p<n^6+6n^5<(n+1)^6.
\]
证明将整数下段、有限实变量区间及无限尾部衔接起来。有限解析段采用非负梯形权：积分显式公式把加权素数和化为实际零点和，密度网格与完整高度余项控制后者；低秩 Taylor 分解把所需超越不等式化为有限有理包围和整数矩证书。两个基准区间的证书经连续平移定理覆盖上中段与远段，有限高度 RH 与缩放核比较补足下中段。无限尾部使用两个端点共享同一截断高度的显式公式，整数下段使用素性见证与修正的显式短区间表。本文按最终 Lean 证明的依赖顺序给出完整数学叙述，列明七项输入的来源、证据和使用位置，并说明全部项目源码、数值证书和重建记录的对应关系。
\end{abstract}
\noindent\textbf{关键词：}素数短区间；准黎曼条件；显式零点密度；平滑显式公式；严格数值证书；Lean。

\medskip
\noindent\textbf{English abstract.}
Under the quasi-Riemann bound $\Re\rho\le7/8$, a pinned explicit zero-density estimate, and five individually stated literature inputs, we prove that every positive integer $n$ admits a prime strictly between $n^6$ and $n^6+6n^5$. A trapezoidal smoothing argument, certified integer moments, continuous parameter translation, finite-height RH, and a shared-truncation explicit formula cover all ranges. The manuscript follows the implemented downstream Lean proof and records the precise external inputs, numerical witnesses, and reproducibility evidence.

\medskip
\noindent\textbf{阅读说明。}本文的数学起点为第\ref{sec:inputs}节的七项输入；其后本项目所需的局部结论均展开证明。作者名称依项目约定指 AI 系统及其贡献，人类作者身份暂未填写。本版完成后提交项目发起人审阅，尚未作为新版推送 GitHub。研究过程的完整版本记录见仓库 \repofile{research-history/}；旧论文及原始证书保留原字节。
\tableofcontents
\clearpage
\section{问题、结论与证明路线}\label{sec:introduction}

相邻整数幂之间的素数问题问：固定整数 $k\ge2$，是否对每个正整数 $n$，开区间 $(n^k,(n+1)^k)$ 中都存在素数？$k=2$ 是通常的勒让德猜想。本文处理 $k=6$，并将每一个外部数学输入固定为可核对的精确陈述。

这一问题有两个尺度。对大 $n$，区间长度约为 $kn^{k-1}$，可以借助零点分布与显式公式控制素数计数；对小 $n$，渐近估计常数较大，直接素性见证和显式短区间结果更有效。两部分之间的有限范围必须被连续覆盖，而不能由若干采样点的成功代替。

\begin{theorem}[主定理]\label{thm:main}
接受第\ref{sec:inputs}节的 H1--H7。对每个整数 $n\ge1$，存在素数 $p$，使
\begin{equation}\label{eq:main}
n^6<p<n^6+6n^5<(n+1)^6.
\end{equation}
\end{theorem}

主定理先证明较短区间内的素数存在性。右侧第二个严格不等式完全来自二项式展开；数论工作集中在第一个区间。对于正实数 $x$，记
\[
\mathcal S_6(x):\quad \exists p\text{ prime},\quad x<p<x+6x^{5/6},
\qquad
\mathcal S_4(x):\quad \exists p\text{ prime},\quad x<p<x+4x^{3/4}.
\]
当 $x\ge1$ 时，$4x^{3/4}<6x^{5/6}$，所以 $\mathcal S_4(x)$ 足以推出 $\mathcal S_6(x)$。这一长度比较允许不同范围采用各自有效的参数。

\subsection{研究前提与目标的选择}
本项目选择接受 QRH$(7/8)$，即实际黎曼 $\zeta$ 函数每个非平凡零点的实部至多为 $7/8$。研究动机来自 OpenAI 公开的相关论文与专项形式化材料\cite{OpenAIQRH}；本文取其明确的零点位置结论作为 H1。显式密度输入另列为 H2，其他五项文献结果列为 H3--H7。这样一来，主定理的数学含义由这些陈述本身固定，上游证据的更新可以独立记录。

路线的阶段研究目标仍为 $k=4$，长期目标为 $k=3$。$k=6$ 是当前已经闭合全部范围、并完成下游形式化的成果。本文会解释哪些估计可复用到较小指数，但只把已经闭合的范围写成定理。

\subsection{覆盖结构}
令 $L=\log x$。最终证明按下表组织；两个相邻范围的共同边界都有明确归属。
\begin{center}\small
\begin{tabular}{p{0.26\linewidth}p{0.25\linewidth}p{0.36\linewidth}}
\toprule
范围 & 局部目标 & 主要方法\\
\midrule
整数 $1\le n\le1242317999994$ & $n^6<p<n^6+6n^5$ & 有限素性见证、CHL 十二行和整数阈值\\
$167\le L\le715/4$ & $\mathcal S_6(x)$ & 有限高度 RH、缩放核及反向密度比较\\
$715/4\le L\le400$ & $\mathcal S_6(x)$ & 困难基准证书及连续传播\\
$400\le L\le700$ & $\mathcal S_4(x)$ & 远段梯形证书及连续传播\\
$700\le L$ & $\mathcal S_4(x)$ & 共同截断显式公式与解析尾界\\
\bottomrule
\end{tabular}\end{center}
整数下段与 $L=167$ 处有已证明的重叠。最终拼接只需将 $x=n^6$ 代入相应实变量结论，再作严格的长度比较。

\subsection{本项目的具体工作}
本文所组织的证明包含四项相互连接的工作。首先，用实际 $\zeta$ 零点及解析重数建立加权显式公式，并把素数幂的贡献从正的加权和中扣除。其次，将零点密度估计变成非负有限网格预算，并处理无限高度中所有残余权重。再次，用指数矩的低秩分解产生可由内核检查的整数证书，再证明证书随实参数的连续传播。最后，把这些估计与有限整数段、共同截断尾部组合成全范围结论。

证明的三个代表性完整预算为 $0.939772$、$0.966283$ 和 $0.450022$，均严格小于归一化主项 $1$。它们分别服务于困难基准区间、下中段和远段。各预算的分项及推导在后文给出；这些数字不是未经证明的程序输出前提。

\subsection{阅读路径}
第\ref{sec:inputs}节固定输入，第\ref{sec:preliminaries}节统一数学对象。第\ref{sec:trapezoid}--\ref{sec:coverage}节建立有限解析段；第\ref{sec:tail}节处理无限尾部；第\ref{sec:lower}节处理整数下段；第\ref{sec:assembly}节完成主定理。后两节讨论方法适用性和可复现形式化。关键命题后的形式化对应列出文件与声明名，方便读者在数学证明和源码之间往返。

\leanref{FinalResult}{sixth_powers_from_named_inputs; sixth_power_short_from_named_inputs}

\section{七项精确外部输入}
\label{sec:inputs}

最终 Lean 定理将下表七项命题逐一列为参数。它们约束真实的黎曼
$\zeta$ 函数及素数，而不预先假定本项目三个解析区间的结论。
本库完成的是从这些输入到六次幂结论的下游推导；本节将输入的数学陈述
与其来源证据分开。所称可信性依据是官方发布、发表记录、来源核对或实际
形式化工作的类别，不是正确概率。

\begin{table}[htbp]
\centering\small
\caption{最终定理的七项外部参数。全部保留在定理类型中。}
\label{tab:inputs}
\begin{tabular}{p{0.06\linewidth}p{0.28\linewidth}p{0.55\linewidth}}
\toprule
编号 & 数学输入 & 现有来源证据与边界 \\
\midrule
H1 & QRH$(7/8)$ & 按研究约定接受 OpenAI 官方证明发布；本库未重建上游。 \\
H2 & CS v2 密度界 & 固定预印本的指定十二行；未确认同行评审发表或接受。 \\
H3 & Rosser 计数界 & 已发表结果，经 Dusart 重述核对；保留来源闭矩形计数。 \\
H4 & 积分显式公式 & 已发表公式的精确特化；大学机构库接受稿已核对。 \\
H5 & 有限高度 RH & 已发表的严格计算结果；本项目未重跑上游零点验证。 \\
H6 & CHJII 公式 & 已发表；所用 arXiv 陈述对应通过，期刊定稿比对待补。 \\
H7 & CHL 十二行 & 原文及勘误已发表；修订表对应通过，仍有版本和见证认证待办。 \\
\bottomrule
\end{tabular}
\end{table}

\subsection{共同记号}

写 $\rho=\beta+\ii\gamma$，令
\[
\mathcal Z=\{\rho\in\C:\zeta(\rho)=0,\ 0<\Re\rho<1\},\qquad
\mathcal Z(T)=\{\rho\in\mathcal Z:|\Im\rho|\le T\}.
\]
每个不同零点以真实解析重数 $m(\rho)$ 计入。
截断参数取 $T\in\R$；当 $T<0$ 时截断集合为空，相关计数取零。
Lean 中的对象分别为
\leanref{ZetaZeros}{NontrivialZero}、
\leanref{ZetaZeros}{allZetaZeros} 和
\leanref{ZetaZeros}{zetaMultiplicity}。
定义
\begin{equation}
\label{eq:zero-counts}
N_+(T)=\sum_{\substack{\rho\in\mathcal Z\\0<\gamma\le T}}m(\rho),
\qquad
N(\sigma,T)=\sum_{\substack{\rho\in\mathcal Z\\\sigma\le\beta,\ 0<\gamma\le T}}m(\rho).
\end{equation}
这里高度上端与实部下端均含等号。$\mathcal Z(T)$ 在定义时保留零高度；
实区间 $(0,1)$ 无零点是本库的另一个已证引理
\leanref{ZetaZeros}{riemannZeta_ne_zero_real_strip}，没有暗中加入输入定义。
所有对数都是自然对数。对正实数 $x$，复幂取
$x^s=\exp(s\log x)$，与 Lean 使用的主支复幂一致。设
\[
\psi(x)=\sum_{1\le n\le\lfloor x\rfloor}\Lambda(n),\qquad
H=3\cdot10^{12};
\]
$\Lambda$ 是通常的 von Mangoldt 函数，$\psi$ 的整数端点包含整项。

\subsection{H1：QRH$(7/8)$}

\begin{definition}[H1：八分之七零点界]
\label{hyp:qrh}
对每个 $\rho\in\mathcal Z$，都有 $\Re\rho\le7/8$。
\end{definition}
此处对应 \leanref{DensityInput}{QRHSevenEighths}。
采用 H1 的研究依据是 OpenAI 于 2026 年 9 月 30 日公布的证明稿
\cite{OpenAIQRH}，固定仓库提交为
\nolinkurl{4eaee87548250b2682affda7d05d947c05554bae}。
该提交的 Family 003 说明报告了 $\zeta$ 的 $7/8$ 无零半平面形式化；
本项目只取其 $\zeta$ 部分，并未导入或独立重建该上游工程。
这是按项目约定接受的无限域输入，不能用有限高度数值扫描替代，也不等同于 RH。

\subsection{H2：固定版本 CS 密度界}

\begin{definition}[H2：CS v2 的指定十二行]
\label{hyp:cs}
令
\[
a(\sigma)=\frac{3(1-\sigma)}{2-\sigma},\qquad
b(\sigma)=\frac{7-5\sigma}{2-\sigma}=a(\sigma)+2.
\]
对表~\ref{tab:cs-input} 的任一行 $j$、所有
$j/32\le\sigma\le(j+1)/32$ 以及所有实数 $T\ge H$，要求
\begin{equation}
\label{eq:cs-input}
N(\sigma,T)\le
B_1 T^{a(\sigma)}(\log T)^{b(\sigma)}
+B_2(\log T)^2+B_3\log T.
\end{equation}
\end{definition}

\begin{table}[htbp]
\centering\small
\caption{H2 的全部系数；有限小数作为精确有理数。}
\label{tab:cs-input}
\begin{tabular}{ccrrr}
\toprule
来源行 $j$ & $\sigma$ 的闭区间 & $B_1$ & $B_2$ & $B_3$ \\
\midrule
16 & $[16/32,17/32]$ & 5.069 & 9.461 & 167.8 \\
17 & $[17/32,18/32]$ & 6.046 & 9.168 & 162.6 \\
18 & $[18/32,19/32]$ & 7.077 & 8.875 & 157.4 \\
19 & $[19/32,20/32]$ & 8.185 & 8.582 & 152.2 \\
20 & $[20/32,21/32]$ & 9.387 & 8.290 & 147.0 \\
21 & $[21/32,22/32]$ & 10.70 & 7.997 & 141.9 \\
22 & $[22/32,23/32]$ & 12.11 & 7.704 & 136.7 \\
23 & $[23/32,24/32]$ & 13.66 & 7.411 & 131.5 \\
24 & $[24/32,25/32]$ & 15.35 & 7.118 & 126.3 \\
25 & $[25/32,26/32]$ & 17.21 & 6.826 & 121.1 \\
26 & $[26/32,27/32]$ & 19.31 & 6.533 & 115.9 \\
27 & $[27/32,28/32]$ & 21.77 & 6.240 & 110.7 \\
\bottomrule
\end{tabular}
\end{table}

该命题是 \leanref{DensityInput}{CSv2}；Lean 行号为
$i=j-16\in\{0,\ldots,11\}$。所用来源为 Chourasiya--Simoni\v{c}
的固定 arXiv v2（2025 年 9 月 30 日），Corollary 1 和 Table 1
\cite{CSv2}。本库的开条带计数是来源计数的限制，因重数非负，限制只减小左端。
现有核查尚未确认该精确输入已获同行评审发表或接受，所以 H2 始终是额外的
预印本前提。其推导与常数来源的独立审查属于另项工作，本稿不代替该审查。
原式中的 $b(\sigma)$ 不可改成 v1 的统一对数三次幂；尾段采用的较粗
$(\log T)^3$ 界由 \leanref{DensityInput}{cs_tail_density}
在明确子域内从 H2 推出。下游数值认证也不能代替 H2 的上游证明。

\subsection{H3：完整 Rosser 计数界}

为保持所引陈述的计数约定，另定义闭矩形计数
\[
N_R(T)=\sum_{\substack{\zeta(\rho)=0\\0\le\Re\rho\le1,\ 0\le\Im\rho\le T}}m(\rho).
\]
\begin{definition}[H3：Rosser 零点计数界]
\label{hyp:rosser}
对每个实数 $T\ge2$，
\begin{equation}
\label{eq:rosser-input}
\begin{split}
\left|N_R(T)-\left(\frac{T}{2\pi}\log\frac{T}{2\pi}
-\frac{T}{2\pi}+\frac78\right)\right|
\le{}&\frac{137}{1000}\log T\\
&+\frac{443}{1000}\log\log T+\frac{397}{250}.
\end{split}
\end{equation}
\end{definition}
对应 \leanref{LiteratureInputs}{RosserZeroCount}。
Rosser 的原文发表于 1941 年的 \emph{American Journal of Mathematics}
\cite{Rosser1941}；本项目通过 Dusart 的正式重述第 877 页式 (10)--(11)
核对系数 \cite{Dusart2016}，并通过作者 HDR 第 34 页和第 82 页
核对闭矩形约定，未声称逐页重证原始论文。
集合包含给出 $N_+(T)\le N_R(T)$，对应
\leanref{LiteratureInputs}{totalZeroCount_le_rosserCount}。
本库已从完整 H3 推出所需的较弱总计数界，而没有额外接受 CS 的总计数辅助式；
Rosser 解析定理本身仍为外部输入。

\subsection{H4：积分显式公式}

令
\[
\Psi_1(x)=\sum_{1\le n\le\lfloor x\rfloor}\Lambda(n)(x-n),\qquad
P(\rho,x)=\frac{x^{\rho+1}}{\rho(\rho+1)},
\]
以及
\[
V(x)=\sum_{j=1}^{\infty}\frac{x^{1-2j}}{2j(2j-1)}.
\]
这分别对应 \leanref{IntegratedPsi}{integratedPsi}、
\leanref{TrapezoidKernel}{zeroPrimitive} 和
\leanref{IntegratedExplicitFormula}{trivialZeroPrimitiveSum}。

\begin{definition}[H4：高度截断的积分显式公式]
\label{hyp:integrated}
对每个实数 $x>1$，在复数域中有
\begin{equation}
\label{eq:integrated-input}
\begin{split}
\lim_{T\to+\infty}\sum_{\rho\in\mathcal Z(T)}m(\rho)P(\rho,x)
={}&\frac{x^2}{2}-\frac{\zeta'(0)}{\zeta(0)}x
+\frac{\zeta'(-1)}{\zeta(-1)}\\
&-V(x)-\Psi_1(x).
\end{split}
\end{equation}
\end{definition}
对应 \leanref{IntegratedExplicitFormula}{IntegratedExplicitFormula}。
Greni\'e--Molteni 的已发表论文 \cite{GM2016} 给出所需精确恒等式：
取数域为有理数域，结合 (1.3a)、(3.8) 和 Lemma 3.3，得到
式~\eqref{eq:integrated-input}，绝对收敛支持这里的闭高度截断极限。
虽然该文题名含 GRH，此一般恒等式不以 GRH 为前提。
机构库接受稿及 arXiv v5 的关键页已经核对；接受稿留数推导中
$s=-1$ 处的局部笔误已记录，其主公式及本库使用的常数符号正确。
Montgomery--Vaughan 第 420 页提供常数项的交叉核对
\cite{MontgomeryVaughan}。
旧源码注释中的 Hildebrand 课件 \cite{Hildebrand} 保留为历史来源，
并非现有文献依据的全部。本库从 H4 推导四点差分及其加权显式公式，
没有把加权公式或误差预算另列为输入；上游复分析证明仍未在本库重建。

\subsection{H5：有限高度 RH}

\begin{definition}[H5：高度不超过 $H$ 的零点]
\label{hyp:finite-rh}
对每个 $\rho\in\mathcal Z$，
\[
0<\Im\rho\le H\quad\Longrightarrow\quad\Re\rho=\frac12.
\]
\end{definition}
对应 \leanref{FiniteRHInput}{FiniteRH3e12}。
Platt--Trudgian 的正式期刊摘要明确给出闭高度上端
$3\cdot10^{12}$ \cite{PlattTrudgian}；arXiv v1 的 Theorem 1
实际达到 $3{,}000{,}175{,}332{,}800$。
本库只使用较小的闭截断，不使用来源另给出的简单性。
该来源是采用区间算术的已发表计算定理，本项目未重跑其大规模零点计算；
从 H5 化简低高度加权零点和的步骤则已在本库形式化。

\subsection{H6：CHJII 的共同高度截断}

\begin{definition}[H6：短区间显式公式]
\label{hyp:chj}
设 $x,h,T\in\R$ 满足
\[
x>0,\quad\log x\ge40,\quad h\ge0,\qquad
\max\{51,(\log x)^2\}<T<\frac{\sqrt{x}-2}{4}.
\]
令 $y=x+h$。则存在一个实数 $T^*\in[T,2T]$，使
\begin{equation}
\label{eq:chj-input}
\left|\psi(y)-\psi(x)-h+
\sum_{\rho\in\mathcal Z(T^*)}m(\rho)\frac{y^\rho-x^\rho}{\rho}\right|
\le\frac{1447}{500T}\bigl(y\log y+x\log x\bigr).
\end{equation}
绝对值表示复模。
\end{definition}
对应 \leanref{LiteratureInputs}{CHJIIShortInterval}。
来源为 Cully-Hugill--Johnston II \cite{CHJII} 的固定 arXiv v3，
Theorem 1.3 与 Table 1 首行
$\alpha=1/2,\omega=0,M=2.894=1447/500$。
两端使用同一个 $T^*$，误差分母是 $T$，并且 $T$ 的两侧限制均为严格不等式。
来源内部的 $\lambda=0.43$ 不另构成使用条件。
该论文的正式发表身份已核实，所引定理、表格及相关证明在 v2/v3 中相同；
期刊定稿与 arXiv 的直接比对仍待补。
本库的独立标量检查得到 $2.8933994128\ldots<2.894$，并用 Arb
认证了九个相关比较；这不等于完整解析定理的重证。
来源审查另检查了本行的纸面全域归约，本库从 H6 推出的无限尾段结论
则属于下游 Lean 证明。

\subsection{H7：CHL 十二行素数区间}

\begin{definition}[H7：指定 CHL 区间结论]
\label{hyp:chl}
对表~\ref{tab:chl-input} 中的每一对 $(X,\Delta)$，以及每个实数
$y\ge X$，存在自然数素数 $p$ 满足
\begin{equation}
\label{eq:chl-input}
y(1-\Delta^{-1})<p\le y.
\end{equation}
\end{definition}
这正是 \leanref{Definitions}{CHLInterval} 及
\leanref{CHLRange}{CHLRows}，没有额外 QRH 条件。

\begin{table}[htbp]
\centering\small
\caption{H7 的全部参数。来源表的第一列为 $\log X$，此处列出 $X$ 本身。}
\label{tab:chl-input}
\begin{tabular}{lcr}
\toprule
Lean 字段 & $X$ & $\Delta$ \\
\midrule
\texttt{row0} & $4\cdot10^{18}$ & 39097000 \\
\texttt{row43} & $\e^{43}$ & 41816800 \\
\texttt{row46} & $\e^{46}$ & 163940000 \\
\texttt{row50} & $\e^{50}$ & 1061200000 \\
\texttt{row55} & $\e^{55}$ & 10288400000 \\
\texttt{row60} & $\e^{60}$ & 76918400000 \\
\texttt{row75} & $\e^{75}$ & 174043000000 \\
\texttt{row90} & $\e^{90}$ & 184304000000 \\
\texttt{row105} & $\e^{105}$ & 191886000000 \\
\texttt{row120} & $\e^{120}$ & 197917000000 \\
\texttt{row135} & $\e^{135}$ & 202553000000 \\
\texttt{row150} & $\e^{150}$ & 207053000000 \\
\bottomrule
\end{tabular}
\end{table}

CHL 原文及其正式勘误均已发表 \cite{CHL,CHLcorr}。
本库实际采用的是作者固定 arXiv v3（2022 年 9 月 14 日）的 Theorem 1
和修订 Table 2；十二行整数、适用域及左开右闭端点已与 Lean 逐项核对。
正式勘误的出版记录及摘要已确认，但全文表格尚未直接比对。
作者代码重放还暴露了一个不同层次的问题：Arb 的精确十进制区间计算
严格确认，公开的已舍入辅助参数生成的 $\Delta$ 在
$\log X=50,60,75,90,105,135$ 六行略低于所报表值，最大相对不足约
$3.36\cdot10^{-6}$。微调辅助参数后，普通高精度计算仍给出正判据余量；
这只是可行性佐证，尚非完整 CHL 判据的严格独立认证。
因此我们保留 H7 的精确陈述，同时公开版本比对和见证认证的未完成项，
不以表格抄录通过代替上游证明。下游从这些区间到六次幂开区间的转换
已在 Lean 中完成。

\begin{remark}[输入与结论的逻辑边界]
七项输入的发表或可信性证据不会使其自动从形式化定理的参数中消失。
本节没有给出七个外部命题在本库中的证明项；给出这些命题后，
\leanref{FinalResult}{sixth_powers_from_named_inputs}
才产生所有自然数 $n\ge1$ 的六次幂结论。
逐项发表状态、来源链接、数值复现范围与未完成项还以机器可读形式保存在
\texttt{documentation/inputs.json}，便于审查者将数学依赖与证明状态分别核查。
\end{remark}

\section{基本对象、零点约定与算术工具}\label{sec:preliminaries}

\subsection{计数函数与端点}
所有对数均为自然对数。对实数 $y\ge0$，定义
\[
\psi(y)=\sum_{1\le m\le\lfloor y\rfloor}\Lambda(m),\qquad
\theta(y)=\sum_{\substack{p\le y\\p\text{ prime}}}\log p,\qquad
\Psi_1(y)=\sum_{1\le m\le\lfloor y\rfloor}\Lambda(m)(y-m).
\]
von Mangoldt 函数满足 $\Lambda(p^a)=\log p$（$a\ge1$），对其余正整数为零。因而 $\psi$ 的增量计入素数和高次素数幂，$\theta$ 的增量只计素数。两者都采用右连续的 $m\le y$ 约定。由于 $\Psi_1$ 在整数端点处的新增项权为零，其定义也等于 $\int_0^y\psi(t)\dd t$，没有半权端点差异。

\subsection{真实零点、重数与有限截断}
本文的非平凡零点是满足 $\zeta(\rho)=0$、$0<\Rea\rho<1$ 的复数。写 $\rho=\beta+\ii\gamma$，以 $m(\rho)$ 记解析零点阶数。所有零点和都带重数，且不假定零点简单。对任意实数 $T$，记
\[
\mathcal Z_T=\{\rho:0<\beta<1,\ |\gamma|\le T,\ \zeta(\rho)=0\},
\quad
\mathcal Z_T^+=\{\rho\in\mathcal Z_T:0<\gamma\le T\}.
\]
这些集合有限：在有界矩形中，亚纯函数的零点是离散的；$s=1$ 的极点附近也不产生零点的聚集。Lean 使用 Mathlib 的紧集零点有限性结论，把实际集合转成有限集合后求和。计数函数为
\[
N(T)=\sum_{\rho\in\mathcal Z_T^+}m(\rho),\qquad
N(\sigma,T)=\sum_{\substack{\rho\in\mathcal Z_T^+\\\beta\ge\sigma}}m(\rho).
\]
这里 $N(T)=N_+(T)$，与第~\ref{sec:inputs} 节的记号相同。
若将截断定义延伸到 $T<0$，相应集合为空、计数为零；本文的解析估计均在各自注明的非负高度范围使用。
增大 $T$ 只添加非负重数项；增大 $\sigma$ 只删除这些项。这给出随后分层求和使用的单调性。

共轭恒等式 $\zeta(\overline s)=\overline{\zeta(s)}$ 同时保持零点与重数。为把两侧高度的和写成正高度和的两倍，还要排除实轴上的非平凡零点。这一步在本项目中有独立证明，不额外列为第八项输入。

\begin{lemma}[实轴无非平凡零点]\label{pre:real-axis}
对 $0<s<1$，有 $\zeta(s)\ne0$。
\end{lemma}
\begin{proof}
给出与形式化实现一致的 theta 核证明。设
\[
\Theta(t)=\sum_{n\in\mathbb Z}\e^{-\pi n^2t},\qquad t>0.
\]
由 theta 变换 $\Theta(t)=t^{-1/2}\Theta(1/t)$，经去除主项的核可取为
\[
K(t)=
\begin{cases}
t^{-1/2}(\Theta(1/t)-1),&0<t<1,\\
\Theta(t)-1,&t>1.
\end{cases}
\]
在 $t=1$ 的定义不影响积分。每项均非负，而且在两端有足够衰减。令
\[
F(s)=\frac12\int_0^\infty t^{s/2-1}K(t)\dd t.
\]
Mathlib 中的完成 zeta 函数给出已形式化的 Mellin 恒等式
\[
\pi^{-s/2}\Gamma(s/2)\zeta(s)
=F(s)-\frac1s-\frac1{1-s},\qquad F(1-s)=F(s).
\]
对 $0\le s\le1$，当 $t\ge1$ 时 $t^{s/2-1}\le1$；当 $0<t<1$ 时 $t^{s/2-1}\le t^{-3/2}$。乘以非负的 $K(t)$ 并积分得
\[
F(s)\le F(2)+F(-1)=2F(2).
\]
由 $\zeta(2)=\pi^2/6$ 和 $\Gamma(1)=1$，有 $F(2)=\pi/6-1/2$。若 $\zeta(s)=0$，上述恒等式将迫使
\[
4\le\frac1{s(1-s)}=\frac1s+\frac1{1-s}=F(s)
\le\frac\pi3-1<\frac13,
\]
矛盾。这里 $s(1-s)\le1/4$，且 $\pi<4$，所有分母均正。
\end{proof}
\leanref{ZetaZeros}{riemannZeta_ne_zero_real_strip; zetaMultiplicity_conj}

因此，对共轭不变的实函数 $f$，
\begin{equation}\label{eq:positive-pair}
\sum_{\rho\in\mathcal Z_T}m(\rho)f(\rho)
=2\sum_{\rho\in\mathcal Z_T^+}m(\rho)f(\rho).
\end{equation}
这是有限集合的配对恒等式，不要求先讨论条件收敛的无限和。
\leanref{ZetaZeros}{sum_allZetaZeros_eq_twice; all_zero_multiplicities_eq_twice_total}

\subsection{正实底数的复幂}
对 $t>0$，统一规定 $t^z=\exp(z\log t)$。其模为 $|t^{\beta+\ii\gamma}|=t^\beta$，导数为 $zt^{z-1}$。若 $x>0$、$h\ge0$、$0<\beta<1$ 且 $\rho\ne0$，积分导数后取模得
\begin{equation}\label{eq:cpow-short}
\left|\frac{(x+h)^\rho-x^\rho}{\rho}\right|
\le\int_x^{x+h}t^{\beta-1}\dd t
\le h x^{\beta-1}.
\end{equation}
最后一步利用 $\beta-1<0$。将它乘以自然数重数并求有限和即可用于截断公式，不需要交换无限求和与积分。
\leanref{CpowInterval}{norm_cpow_sub_div_le}

\subsection{素数幂贡献与严格开区间}
对 $Y\ge1$，Mathlib 已证明无条件估计
\begin{equation}\label{eq:psi-theta-sqrt}
0\le\psi(Y)-\theta(Y)\le2\sqrt Y\log Y.
\end{equation}
若 $S\subseteq\{1,\ldots,\lfloor Y\rfloor\}$，权函数满足 $w(m)\le W$、$W\ge0$，则由 $\Lambda(m)\ge0$，
\[
\sum_{\substack{m\in S\\m\text{ 非素数}}}\Lambda(m)w(m)
\le W(\psi(Y)-\theta(Y))\le2W\sqrt Y\log Y.
\]
此估计不要求事先把所有素数幂逐一枚举。当 $\log Y\ge4\log2$ 时，也得到方便与原始证书对齐的较弱形式
\begin{equation}\label{eq:prime-power-log-square}
\sum_{\substack{m\in S\\m\text{ 非素数}}}\Lambda(m)w(m)
\le\frac{W\sqrt Y(\log Y)^2}{2\log2}.
\end{equation}
\leanref{PrimePowerBudget}{weightedNonprime_le_sqrt_log; weightedNonprime_le_manuscript_budget}

加权证明中，若正权的支持严格位于 $(x,y)$，且扣除非素数贡献后总和仍正，那么有限和中至少有一项来自严格内部的素数。未加权的 $\psi(y)-\psi(x)$ 则含右端点，必须另外留出端点贡献。

\begin{lemma}[未加权的严格端点提取]\label{pre:endpoint}
设 $0\le x\le y$、$y\ge1$、$\log y\ge4\log2$。若
\[
\psi(y)-\psi(x)>
\frac{\sqrt y(\log y)^2}{2\log2}+\log y,
\]
则 $(x,y)$ 内含有素数。
\end{lemma}
\begin{proof}
假设开区间内没有素数。在 $(x,y]$ 中至多可能有右端点 $y$ 自身这一个素数，所以 $\theta(y)-\theta(x)\le\log y$。又因 $\theta(x)\le\psi(x)$，由式\eqref{eq:prime-power-log-square}（$W=1$），
\[
\psi(y)-\psi(x)
\le(\psi(y)-\theta(y))+(\theta(y)-\theta(x))
\le\frac{\sqrt y(\log y)^2}{2\log2}+\log y,
\]
与假设矛盾。
\end{proof}
\leanref{PrimePowerBudget}{theta_sub_le_log_of_no_prime; prime_of_psi_sub_gt_manuscript_budget}

这些准备将后文的解析任务固定为：找到一个严格正的素数贡献下界。梯形权通过零端点支持处理边界，尾部通过额外的 $\log y$ 预算处理边界；二者均得出真正的开区间结论。

\section{梯形权、积分显式公式与素数提取}\label{sec:trapezoid}

本节把素数存在性归结为三个具有明确含义的非负上界之和小于
主项。论证首先是一个有限和恒等式，然后才使用积分显式公式。
这种顺序使端点、真素数幂和零点截断的处理均可单独核对。
本节使用的显式公式是前述输入
\texttt{IntegratedExplicitFormula}；它本身是经典的无条件恒等式，
其精确来源及与本文约定的对应见文献 \cite{GM2016} 和输入审计。
这里不额外假定 RH，也不把 QRH 用于推导该恒等式。

\subsection{从有限 Mangoldt 和构造梯形权}

对 $t>1$ 定义
\begin{equation}\label{tr:psi1}
 \Psi_1(t)=\sum_{1\leq m\leq\lfloor t\rfloor}\Lambda(m)(t-m).
\end{equation}
若 $t$ 是整数，最后一项为零；因此这一约定没有端点半权问题。
对于任何 $Y\geq t$，也可以写成
\[
 \Psi_1(t)=
 \sum_{1\leq m\leq\lfloor Y\rfloor}\Lambda(m)(t-m)_+,
 \qquad v_+=\max\{v,0\}.
\]
因为新增的每一项均为零。这是下面交换四点差分与有限求和的依据。

设 $x>1$、$h>0$，并设 $0<a\leq b$、$a+b=1$。定义
\begin{align}
 D_{x,h}F
  &=\frac{F(x)-F(x+ah)-F(x+bh)+F(x+h)}{abh},
       \label{tr:difference}\\
 w_{x,h}(t)
  &=\frac{(x-t)_+-(x+ah-t)_+-(x+bh-t)_++(x+h-t)_+}
          {abh}. \label{tr:weight}
\end{align}
在不引起混淆时省略 $x,h$ 下标。逐段展开正部函数得到
\begin{equation}\label{tr:piecewise}
 w(t)=
 \begin{cases}
 0,&t\leq x,\\
 (t-x)/(abh),&x\leq t\leq x+ah,\\
 1/b,&x+ah\leq t\leq x+bh,\\
 (x+h-t)/(abh),&x+bh\leq t\leq x+h,\\
 0,&t\geq x+h.
 \end{cases}
\end{equation}
相邻公式在公共端点取相同值。特别地，
\begin{equation}\label{tr:weight-properties}
 0\leq w(t)\leq \frac1b,\qquad
 w(t)>0\Longrightarrow x<t<x+h,\qquad
 \int_{\R}w(t)\,\dd t=h.
\end{equation}
面积归一化为 $h$，故后面的主项也是 $h$。
本文固定
\begin{equation}\label{tr:fixed-ab}
 a=\frac1{32},\qquad b=\frac{31}{32}.
\end{equation}
该选择不表示在所有参数下都最优；它是当前证书实际采用的权。

\begin{lemma}[有限和与四点差分]\label{tr:finite-identity}
设 $Y=x+h$，并令
\[
 W(x,h)=\sum_{1\leq m\leq\lfloor Y\rfloor}\Lambda(m)w_{x,h}(m).
\]
则 $W(x,h)=D_{x,h}\Psi_1$。此外，
$D_{x,h}(t^2/2)=h$，而 $D_{x,h}$ 消去一切仿射函数。
\end{lemma}
\begin{proof}
将 \eqref{tr:psi1} 的四个截断统一扩展到 $Y$，
并把四个有限和合并，其第 $m$ 项恰为
$\Lambda(m)w_{x,h}(m)$。常数项的系数为 $1-1-1+1=0$；
一次项的系数为 $x-(x+ah)-(x+bh)+(x+h)=0$。
二次项的分子为 $abh^2$，除以 $abh$ 即为 $h$。
\end{proof}

这些有限和及支撑性质对应
\leanref{IntegratedPsi}{weightedMangoldt_trapezoid_eq}、
\leanref{IntegratedPsi}{trapezoidWeight_bounds}
和 \leanref{IntegratedPsi}{trapezoidWeight_pos_support}。

\subsection{零点核的三个互补上界}

对 $\rho=\beta+\ii\gamma\notin\{0,-1\}$ 和 $t>0$，令
\[
 G_\rho(t)=\frac{t^{\rho+1}}{\rho(\rho+1)},\qquad
 K_{x,h}(\rho)=D_{x,h}G_\rho,\qquad
 t^\rho=\exp(\rho\log t).
\]
正实轴上的幂使用实对数，因而这里不存在复对数分支选择。
直接求导得 $G_\rho'(t)=t^\rho/\rho$ 及
$G_\rho''(t)=t^{\rho-1}$。

\begin{lemma}[基本界、一次衰减界与二次衰减界]\label{tr:three-kernels}
在上述条件下，若 $\beta\leq1$，则
\begin{equation}\label{tr:base-kernel}
 |K_{x,h}(\rho)|\leq h x^{\beta-1}.
\end{equation}
若另有 $\gamma\ne0$，则
\begin{align}
 |K_{x,h}(\rho)|
 &\leq h x^{\beta-1}\frac{2}{b}
                  \frac{x+h}{h|\gamma|},\label{tr:first-kernel}\\
 |K_{x,h}(\rho)|
 &\leq h x^{\beta-1}\frac{4}{ab}
                  \frac{(x+h)^2}{h^2\gamma^2}. \label{tr:second-kernel}
\end{align}
\end{lemma}
\begin{proof}
两次微积分基本定理给出
\[
 K_{x,h}(\rho)=
 \frac1{abh}\int_0^{ah}\int_0^{bh}
       (x+u+v)^{\rho-1}\,\dd v\,\dd u.
\]
由于 $\beta-1\leq0$，被积函数的模不超过
$x^{\beta-1}$，得到 \eqref{tr:base-kernel}。
Lean 中使用两次导数范数的增量界证明同一估计。

将分子分成两段长度为 $ah$ 的增量。
在 $x\leq t\leq x+h$ 上，
\[
 |G_\rho'(t)|
 =\frac{t^\beta}{|\rho|}
 \leq \frac{(x+h)x^{\beta-1}}{|\gamma|}.
\]
两个增量的模之和除以 $abh$，得到
\eqref{tr:first-kernel}。最后，
$|\rho|\geq|\gamma|$、$|\rho+1|\geq|\gamma|$，且
$t^{\beta+1}=t^2t^{\beta-1}\leq(x+h)^2x^{\beta-1}$。
分别估计四个函数值即得 \eqref{tr:second-kernel}。
\end{proof}

令
\begin{equation}\label{tr:height-kernel}
 c_1=\frac{64}{31},\qquad c_2=\frac{4096}{31},\qquad
 g_q(t)=\min\left\{1,\frac{c_1q}{t},\frac{c_2q^2}{t^2}\right\}
 \quad(t>0,\ q\geq0).
\end{equation}
由 \eqref{tr:fixed-ab}，$c_1=2/b$，$c_2=4/(ab)$。
若 $q\geq(x+h)/h$，三个界合并为
\begin{equation}\label{tr:kernel-min}
 |K_{x,h}(\rho)|\leq h x^{\beta-1}g_q(|\gamma|).
\end{equation}
$g_q$ 关于高度单调不增，关于 $q$ 单调不减，而且非负。
其两个衰减分支具有不同用途：一次分支改善有限高度范围的估计，
二次分支使无限高度余项可被几何级数控制。
对应的形式化结果为
\leanref{TrapezoidKernel}{norm_trapezoidKernel_le_base}、
\leanref{TrapezoidKernel}{norm_trapezoidKernel_le_first}、
\leanref{TrapezoidKernel}{norm_trapezoidKernel_le_second}
及 \leanref{TrapezoidKernel}{norm_trapezoidKernel_le_min}。

\subsection{积分显式公式的四点作用及极限}

记 $m_\rho$ 为非平凡零点的重数。使用按
$|\operatorname{Im}\rho|\leq T$ 截断的零点和，输入公式可写为
\begin{equation}\label{tr:explicit}
 \Psi_1(t)=\frac{t^2}{2}
 -\lim_{T\to+\infty}
   \sum_{|\operatorname{Im}\rho|\leq T}m_\rho G_\rho(t)
 -\frac{\zeta'}{\zeta}(0)t+\frac{\zeta'}{\zeta}(-1)-Q(t),
 \qquad t>1,
\end{equation}
其中
\begin{equation}\label{tr:trivial-primitive}
 Q(t)=\sum_{j=0}^{\infty}
       \frac{t^{-2j-1}}{2(j+1)(2j+1)}
     =\sum_{j=0}^{\infty}G_{-2(j+1)}(t).
\end{equation}
这里的两个仿射常数是精确常数，不能把未指定的 $O(t)$
直接当作可消去的仿射函数。不过，一旦采用精确公式
\eqref{tr:explicit}，其数值完全不进入后续预算，因为四点差分消去它们。

定义
\[
 Z_T(x,h)=\sum_{|\operatorname{Im}\rho|\leq T}
                  m_\rho K_{x,h}(\rho),\qquad E(x,h)=D_{x,h}Q .
\]
对四个大于 $1$ 的实数同时使用 \eqref{tr:explicit}，并对四个极限作有限线性组合，
由引理 \ref{tr:finite-identity} 得
\begin{equation}\label{tr:zero-limit}
 Z_T(x,h)\longrightarrow h-W(x,h)-E(x,h).
\end{equation}
这一步只需要输入已经给出的高度截断极限；无需在四点操作中
另行假定非平凡零点级数绝对收敛。
因此，只要对每个 $T$ 证明 $|Z_T(x,h)|\leq hZ$，
闭球对极限封闭就给出
\begin{equation}\label{tr:weighted-lower}
 W(x,h)\geq h(1-Z)-|E(x,h)|.
\end{equation}
本节以下所有零点预算都对任意截断高度成立，正是为了这一极限步骤。
对应
\leanref{IntegratedExplicitFormula}{trapezoid_zero_sum_tendsto}。

\begin{lemma}[平凡零点的完整上界]\label{tr:trivial-bound}
若 $x>1$，则
\begin{equation}\label{tr:V}
 |E(x,h)|\leq h V(x),\qquad
 V(x)=\frac{x^{-3}}{1-x^{-2}}.
\end{equation}
\end{lemma}
\begin{proof}
对 $t>1$，\eqref{tr:trivial-primitive} 中第 $j$ 项的模
不超过 $t^{-1}(t^{-2})^j$，故其级数绝对收敛。
对 $\rho=-2(j+1)$ 使用基本核界 \eqref{tr:base-kernel} 得
\[
 |K_{x,h}(-2(j+1))|
 \leq h x^{-2j-3}=h x^{-3}(x^{-2})^j.
\]
此级数也绝对收敛。于是四个收敛级数的有限线性组合可以逐项进行，
且几何级数求和给出 \eqref{tr:V}。
\end{proof}
对应 \leanref{TrapezoidTrivial}{norm_trapezoidTrivialDifference_le}。

\subsection{真素数幂的显式预算}

记 $\theta(Y)=\sum_{p\leq Y}\log p$、
$\psi(Y)=\sum_{m\leq Y}\Lambda(m)$。由 $0\leq w\leq1/b$，
非素数项的总贡献满足
\begin{equation}\label{tr:nonprime}
 \sum_{\substack{1\leq m\leq Y\\m\ \text{非素数}}}
       \Lambda(m)w(m)
 \leq\frac1b\bigl(\psi(Y)-\theta(Y)\bigr)
 \leq\frac{\sqrt Y(\log Y)^2}{2b\log2}.
\end{equation}
后一界可以直接由真素数幂计数看出：
若 $p^\ell\leq Y$ 且 $\ell\geq2$，则
$\ell\leq\log Y/\log2$，每个指数对应的底数数量不超过 $\sqrt Y$，
而 $\log p\leq(\log Y)/2$；对这些可能性求和只会多计。
在本文使用范围内，Lean 另由
$\psi(Y)-\theta(Y)\leq2\sqrt Y\log Y$ 和 $\log Y\geq4\log2$
推出同一个较宽上界。
两种推导都无须假定零点分布，也无须判定区间内有哪些素数幂。
见 \leanref{PrimePowerBudget}{weightedNonprime_le_psi_sub_theta}
及 \leanref{PrimePowerBudget}{weightedNonprime_le_manuscript_budget}。

为了给出完全确定的单元预算，令 $k\in\{4,6\}$，
\begin{equation}\label{tr:domains}
 x=\e^L,\qquad h_k(L)=k\e^{(1-1/k)L},\qquad
 \epsilon_k(A)=k\e^{-A/k},\qquad
 q_k(L)=1+\frac{\e^{L/k}}{k}.
\end{equation}
于是 $x+h_k(L)=\e^L(1+\epsilon_k(L))$，且
$q_k(L)=(x+h_k(L))/h_k(L)$。对 $0<A\leq L\leq C$ 定义
\begin{equation}\label{tr:Pcell}
 P_k(A,C)=
 \frac{32}{31}\,
 \frac{\sqrt{1+\epsilon_k(A)}
       \bigl[C+\log(1+\epsilon_k(A))\bigr]^2}
      {2k\log2}\,
 \e^{-(1/2-1/k)A},
 \qquad
 V_A=\frac{\e^{-3A}}{1-\e^{-2A}}.
\end{equation}
这些是后文 $P$ 和 $V$ 的完整定义。特别是
$P_6$ 的分母为 $12\log2$、指数为 $-A/3$；
$P_4$ 的分母为 $8\log2$、指数为 $-A/4$。

\begin{lemma}[单元内的统一小项界]\label{tr:cell-small}
在本文的中段范围 $167\leq A\leq L\leq C$（$k=6$）
或远段范围 $400\leq A\leq L\leq C$（$k=4$）内，
非素数加权和不超过 $h_k(L)P_k(A,C)$，
且 $|E(\e^L,h_k(L))|\leq h_k(L)V_A$。
\end{lemma}
\begin{proof}
把 \eqref{tr:nonprime} 除以 $h_k(L)$，使用
\[
 \frac{\sqrt{x+h_k(L)}}{h_k(L)}
 =\frac1k\sqrt{1+\epsilon_k(L)}
           \e^{-(1/2-1/k)L}.
\]
函数 $\epsilon_k$ 递减，且 $1/2-1/k>0$。
又有
$\log(x+h_k(L))\leq C+\log(1+\epsilon_k(A))$。
所有因子非负，逐因子放大即得 $P_k(A,C)$。
对平凡零点，$\e^{-3L}\leq\e^{-3A}$，
$1-\e^{-2L}\geq1-\e^{-2A}>0$，使用 \eqref{tr:V} 即可。
\end{proof}
此处分别对应
\leanref{MiddleDomains}{middle_prime_power_normalized_le}、
\leanref{MiddleDomains}{middle_prime_power_le}
和 \leanref{FarDomains}{far_prime_power_le}。

\subsection{严格不等式如何产生真正的素数}

\begin{proposition}[三项预算判据]\label{tr:prime-criterion}
设输入积分显式公式成立，$x>1$、$h>0$。
若对任意实数 $T$ 均有 $|Z_T(x,h)|\leq hZ$，
平凡零点差分满足 $|E(x,h)|\leq hV$，
非素数加权和不超过 $hP$，且
\begin{equation}\label{tr:criterion}
 Z+P+V<1,
\end{equation}
则存在素数 $p$ 满足 $x<p<x+h$。
\end{proposition}
\begin{proof}
由 \eqref{tr:weighted-lower}，
$W(x,h)\geq h(1-Z-V)>hP$。
减去不超过 $hP$ 的非素数项，剩下的素数加权和严格为正。
这是有限个非负项的和，所以至少有一项
$\log p\,w(p)>0$。由 \eqref{tr:weight-properties} 的支撑性质，
$x<p<x+h$。由于两个端点权均为零，结论是严格开区间，
不存在把端点素数误计入结论的问题。
\end{proof}
该结论对应 \leanref{TrapezoidPrime}{prime_of_trapezoid_budgets}。
后续任务因此明确为：构造完整零点预算 $Z$，再与
\eqref{tr:Pcell} 中已定义的 $P_k(A,C)$、$V_A$ 相加。
本节没有把任何待证的素数存在命题作为输入。

\section{零点密度网格与完整高度余项}\label{sec:density}

本节从 QRH$(7/8)$、CS v2 和 Rosser 零点计数输入推出
第 \ref{sec:trapezoid} 节所需的统一零点预算。
核心是两次正权离散化：先沿零点实部方向使用密度计数，
再沿高度方向使用递减核。每一次离散化的误差方向都是向上；
高度网格之外的贡献由一个明确的几何级数估计，不能省略。

\subsection{重数约定与 CS v2 的精确使用方式}

只对非平凡零点 $\rho=\beta+\ii\gamma$ 计数，并始终乘重数
$m_\rho$。令
\[
 N(T)=\sum_{0<\gamma\leq T}m_\rho,\qquad
 N(\sigma,T)=\sum_{\substack{0<\gamma\leq T\\\beta\geq\sigma}}m_\rho,
 \qquad H=3\cdot10^{12}.
\]
这里 $N(T)$ 与前文的 $N_+(T)$ 是同一计数的简写，
不是包含正、负高度的总和；$N$ 的第一参数省略时表示全部正高度零点。
闭条件 $\gamma\leq T$ 与 $\beta\geq\sigma$ 均与当前 Lean 输入一致。
若文献采用一个更大的零点集合，限制到本文集合只会减小计数。
共轭对称和重数保持允许把正、负高度的范数和合并为正高度的两倍。
非平凡实零点不存在，故没有 $\gamma=0$ 的剩余项。

对 $r=0,\ldots,11$，CS v2 输入 \cite{CSv2} 在
\[
 \frac{16+r}{32}\leq\sigma\leq\frac{17+r}{32},\qquad T\geq H
\]
上给出
\begin{equation}\label{dens:CS}
 N(\sigma,T)\leq
 B_{1,r}T^{p(\sigma)}(\log T)^{p(\sigma)+2}
 +B_{2,r}(\log T)^2+B_{3,r}\log T,
 \qquad p(\sigma)=\frac{3(1-\sigma)}{2-\sigma}.
\end{equation}
表 \ref{dens:constants} 的有限小数都按所写的精确有理数使用，
不是再次拟合所得的浮点常数。本文承认此条预印本输入；
本节不证明它的上游推导。

\begin{table}[htbp]
\centering
\caption{当前证明实际使用的 CS v2 前十二行常数。}\label{dens:constants}
\begin{tabular}{c|rrr@{\qquad}c|rrr}
$r$&$B_{1,r}$&$B_{2,r}$&$B_{3,r}$
 &$r$&$B_{1,r}$&$B_{2,r}$&$B_{3,r}$\\ \hline
0&5.069&9.461&167.8&6&12.11&7.704&136.7\\
1&6.046&9.168&162.6&7&13.66&7.411&131.5\\
2&7.077&8.875&157.4&8&15.35&7.118&126.3\\
3&8.185&8.582&152.2&9&17.21&6.826&121.1\\
4&9.387&8.290&147.0&10&19.31&6.533&115.9\\
5&10.70&7.997&141.9&11&21.77&6.240&110.7
\end{tabular}
\end{table}

式 \eqref{dens:CS} 中对数的指数是 $p(\sigma)+2$，
其恒等式形式为 $(7-5\sigma)/(2-\sigma)$。
在有限网格证书里不能把它任意换成旧版本的指数；
后文若使用 $\log^3T$，只能作为由
$p(\sigma)+2\leq3$ 和 $\log T\geq1$ 推出的较粗上界。
精确输入和这个代数等式分别见
\leanref{DensityInput}{CSv2} 和
\leanref{DensityInput}{cs_log_exponent_eq}。

Rosser 输入在本文所用高度上蕴含简便估计
\begin{equation}\label{dens:rosser-simple}
 N(\e^u)\leq\e^u u\qquad(\e^u\geq H,\ u\geq1).
\end{equation}
具体地，记源文闭矩形计数为 $N_R(T)$，则 $N(T)\leq N_R(T)$。
由 Rosser 上界、$\log T\geq1$、$\log\log T\leq\log T$ 和
$\pi\geq1$，丢弃负项后有
\[
 N(T)\leq\frac T2\log T+0.580\log T+2.463
       \leq\left(\frac T2+4\right)\log T
       \leq T\log T.
\]
最后一步只需 $T\geq8$，在本文高度域内自动满足。
这里未把 $N(T)\leq T\log T$ 作为第八个外部输入；
对应 \leanref{LiteratureInputs}{totalZeroCount_le_mul_log}。
原始零点计数来源为 \cite{Rosser1941,Dusart2016}。

\subsection{实部网格：从单个指数到密度计数}

设
\[
 \frac12=\sigma_0\leq\sigma_1\leq\cdots\leq\sigma_J=\frac78
\]
是有限网格。每个左端点 $\sigma_j$ 指定一个适用的 CS 行
$r(j)$，即 $\sigma_j$ 属于该行的闭区间。
对 $A\geq0$ 定义
\begin{equation}\label{dens:weights}
 W_j(A)=\e^{(\sigma_{j+1}-1)A}
             -\e^{(\sigma_j-1)A}\geq0.
\end{equation}

\begin{lemma}[向上的阶梯界]\label{dens:staircase}
若 $\beta\leq7/8$，则
\begin{equation}\label{dens:staircase-eq}
 \e^{(\beta-1)A}\leq
 \e^{-A/2}+
 \sum_{j=0}^{J-1}\mathbf{1}_{\{\sigma_j\leq\beta\}}W_j(A).
\end{equation}
\end{lemma}
\begin{proof}
若 $\beta<\sigma_0$，右边至少为 $\e^{-A/2}$，结论由单调性得到。
否则，取最后一个满足 $\sigma_j\leq\beta$ 的网格位置。
右边的增量和望远镜相消，等于一个不小于
$\e^{(\beta-1)A}$ 的右端点值。若 $\beta=\sigma_J$，
所有增量都被计入，右边恰等于 $\e^{(\sigma_J-1)A}$。
在内部网格端点使用闭条件可能多计一个增量，但仍为上界。
\end{proof}

定义正高度的带权累积计数
\[
 F_A(T)=\sum_{0<\gamma\leq T}m_\rho\e^{(\beta-1)A}.
\]
QRH$(7/8)$ 使引理 \ref{dens:staircase} 可用于每一个零点。
将其乘以重数并求和，得到
\begin{equation}\label{dens:weighted-count}
 F_A(T)\leq\e^{-A/2}N(T)+
              \sum_{j=0}^{J-1}W_j(A)N(\sigma_j,T).
\end{equation}
这里无需零点恰好落在某个半开网格单元内；
使用累积计数已处理所有边界重合。

把 $T=\e^u$ 代入，记 $p_j=p(\sigma_j)$，并定义
\begin{align}
 C_j(u)&=
 B_{1,r(j)}\e^{p_ju}u^{p_j+2}
       +B_{2,r(j)}u^2+B_{3,r(j)}u,\label{dens:C}\\
 \Phi_A(u)&=\e^{-A/2+u}u+
                     \sum_{j=0}^{J-1}W_j(A)C_j(u).
                                                   \label{dens:Phi}
\end{align}
当 $\e^u\geq H$、$u\geq1$ 时，
\begin{equation}\label{dens:F-Phi}
 F_A(\e^u)\leq\Phi_A(u).
\end{equation}
至此，所有零点位置都已被一个显式的、有限项的实函数替换。
对应 \leanref{FiniteDensityGrid}{exp_le_grid_staircase}、
\leanref{FiniteDensityGrid}{positiveZeroWeight_le_density_grid}
和 \leanref{FiniteDensityGrid}{positiveZeroWeight_le_gridPhi}。

\subsection{密度上界沿高度方向的增长率}

为了控制无限余项，先证明一个不依赖任何数值采样的增长界。
\begin{lemma}\label{dens:growth}
若 $U\geq6$、$v\geq0$，则
\begin{equation}\label{dens:growth-eq}
 \Phi_A(U+v)\leq\e^{3v/2}\Phi_A(U).
\end{equation}
\end{lemma}
\begin{proof}
对 $u>0$、$v\geq0$、$r\geq0$，由
$\log(1+v/u)\leq v/u$ 得
\[
 \e^{p(u+v)}(u+v)^r
 \leq \e^{(p+r/u)v}\e^{pu}u^r.
\]
在 $1/2\leq\sigma_j\leq7/8$ 上，$0\leq p_j\leq1$，
且 $2\leq p_j+2\leq3$。
因此 $p_j+(p_j+2)/U\leq1+3/6=3/2$。
同样的估计也适用于 $u^2$、$u$ 和基线 $\e^u u$：
相应的 $(p,r)$ 分别为 $(0,2)$、$(0,1)$、$(1,1)$。
式 \eqref{dens:Phi} 中的系数均非负，逐项相加即得结论。
\end{proof}
特别地，将高度乘以 $4$ 时，对数高度增加 $\log4$，
从而 $\Phi_A(U+\log4)\leq8\Phi_A(U)$。
该界对应 \leanref{DensityGridGrowth}{densityGridPhi_shift_le}。

\subsection{高度网格及其不能遗漏的剩余权}

选择单调对数高度网格
$u_0\leq u_1\leq\cdots\leq u_M=U$，并要求
\begin{equation}\label{dens:height-conditions}
 \e^{u_0}\geq H,\qquad u_0\geq1,\qquad U\geq6,\qquad
 g_q(\e^{u_0})=1.
\end{equation}
记 $t_i=\e^{u_i}$、$R=\e^U$，并定义非负质量
\begin{equation}\label{dens:mu}
 \mu_i=g_q(t_i)-g_q(t_{i+1})\geq0,\qquad 0\leq i<M.
\end{equation}
对于任意 $t>0$，有
\begin{equation}\label{dens:height-staircase}
 g_q(t)\leq
 \sum_{i=0}^{M-1}\mu_i\mathbf{1}_{\{t\leq t_{i+1}\}}
        +g_q(\max\{t,R\}).
\end{equation}
若 $t\leq t_0$，右边望远镜相消为 $1$；
若 $t_i<t\leq t_{i+1}$，右边不小于 $g_q(t_i)$；
若 $t>R$，右边等于 $g_q(t)$。
这证明了 \eqref{dens:height-staircase}，并说明余项具有明确的含义：
\emph{即使一个零点位于 $R$ 以下，它在有限高度和之外仍留下
$g_q(R)$ 这一份权。}
因此，仅估计 $\gamma>R$ 的零点尾和不足以完成本证明。

对任意有限截断 $T$，将 \eqref{dens:height-staircase}
乘以 $m_\rho\e^{(\beta-1)A}$ 并对正高度零点求和，
有限网格部分不超过
\[
 \sum_{i=0}^{M-1}\mu_i F_A(t_{i+1})
 \leq\sum_{i=0}^{M-1}\mu_i\Phi_A(u_{i+1}).
\]
剩余部分使用
\begin{equation}\label{dens:quadratic-residual}
 g_q(\max\{t,R\})\leq
             \frac{c_2q^2}{\max\{t,R\}^2}.
\end{equation}

\begin{lemma}[完整二次余项]\label{dens:residual}
在 \eqref{dens:height-conditions} 的条件下，对每个 $T$，
\begin{equation}\label{dens:residual-eq}
 \sum_{0<\gamma\leq T}
 \frac{m_\rho\e^{(\beta-1)A}c_2q^2}{\max\{\gamma,R\}^2}
 \leq16c_2q^2\e^{-2U}\Phi_A(U).
\end{equation}
\end{lemma}
\begin{proof}
若截断集合为空，结论显然。否则选取整数 $N\geq0$，
使集合内所有高度均不超过 $R4^{N+1}$。
对于这些高度，逐项有
\[
 \frac{c_2q^2}{\max\{t,R\}^2}
 \leq\frac{c_2q^2}{R^2}
       \sum_{k=0}^N16^{-k}\mathbf{1}_{\{t\leq R4^{k+1}\}}.
\]
若 $t\leq R$，右边的 $k=0$ 项已足够；
若 $R4^k<t\leq R4^{k+1}$，右边第 $k$ 项已足够。
对零点求和，并用 \eqref{dens:F-Phi} 与引理
\ref{dens:growth}，得到
\begin{align*}
 \text{左边}
 &\leq \frac{c_2q^2}{R^2}
   \sum_{k=0}^N16^{-k}\Phi_A(U+(k+1)\log4)\\
 &\leq \frac{c_2q^2}{R^2}\Phi_A(U)
   \sum_{k=0}^N16^{-k}8^{k+1}\\
 &\leq \frac{c_2q^2}{R^2}\Phi_A(U)\,
          8\sum_{k=0}^{\infty}2^{-k}
   =16c_2q^2\e^{-2U}\Phi_A(U).
\end{align*}
得到的常数与 $T$、$N$ 均无关。
\end{proof}
这给出常数 $16$ 的完整来源。它已经包含高度 $R$ 以下的剩余权，
以及 $R$ 以上全部高度；后续不会再另留一个未定义的尾项。
对应 \leanref{GeometricHeightTail}{finite_quadratic_residual_bound}
和 \leanref{TrapezoidHeightBound}{positive_zero_quadratic_residual_le}。

\subsection{从密度预算回到梯形零点和}

定义可供数值认证的完整正高度预算
\begin{equation}\label{dens:B}
 \mathcal B(A,q;\boldsymbol u)=
 \sum_{i=0}^{M-1}
   \bigl[g_q(\e^{u_i})-g_q(\e^{u_{i+1}})\bigr]\Phi_A(u_{i+1})
       +16c_2q^2\e^{-2U}\Phi_A(U).
\end{equation}
前面的论证表明，对任意 $T$，
\begin{equation}\label{dens:positive-final}
 \sum_{0<\gamma\leq T}m_\rho\e^{(\beta-1)A}g_q(\gamma)
 \leq\mathcal B(A,q;\boldsymbol u).
\end{equation}

\begin{proposition}[完整零点预算]\label{dens:zero-budget}
设 QRH$(7/8)$、CS v2 和 Rosser 输入成立。
若 $x=\e^L$、$0\leq A\leq L$、$h>0$、
$q\geq(x+h)/h$，网格满足以上条件，则对任意 $T$，
\begin{equation}\label{dens:zero-budget-eq}
 |Z_T(x,h)|\leq2h\,\mathcal B(A,q;\boldsymbol u).
\end{equation}
所以第 \ref{sec:trapezoid} 节的判据可取
$Z=2\mathcal B(A,q;\boldsymbol u)$。
\end{proposition}
\begin{proof}
由 QRH，$\beta\leq7/8<1$，故 $L\geq A$ 蕴含
$\e^{(\beta-1)L}\leq\e^{(\beta-1)A}$。
用引理 \ref{tr:three-kernels} 的合并界、三角不等式和共轭配对，
\[
 |Z_T(x,h)|
 \leq2h\sum_{0<\gamma\leq T}
                  m_\rho\e^{(\beta-1)A}g_q(\gamma).
\]
再用 \eqref{dens:positive-final} 即可。
\end{proof}
对应 \leanref{TrapezoidHeightBound}{positive_zero_envelope_le_heightBudget}
和
\leanref{TrapezoidHeightBound}{norm_truncatedTrapezoidZeroSum_le_heightBudget}。
结合命题 \ref{tr:prime-criterion}，每个待认证单元的充分条件现为
\[
 2\mathcal B(A,q;\boldsymbol u)+P_k(A,C)+V_A<1.
\]
下一节把这里的有限实数表达式转为可审阅、可逐项重放的有理证书。

\section{低秩 Taylor 矩、有理包围与两个试点证书}\label{sec:certificates}

本节证明两个具体单元的预算严格小于 $1$。
计算的数学对象是式 \eqref{dens:B} 给出的完整预算，
包括无限高度剩余权；数值证书不替代前两节的解析推导。
我们先把双重有限和分解为非负矩，再将每个实数因子向上包围为
分母 $2^{48}$ 的有理数，最后仅用整数与有理数运算验证目标界。
生成器可以提出这些有理数，但每一个包围关系另有 Lean 证明。

\subsection{把完整零点预算精确拆成主项与辅助项}

沿用实部网格 $\sigma_0,\ldots,\sigma_J$、高度网格
$u_0,\ldots,u_M=U$、左端点 $A$ 和参数 $q$，并假定 $M\geq1$。
这一条件保证完整高度余项可以放入最后一个高度单元。
对 $0\leq i<M$ 记
\begin{align}
 s_i&=u_{i+1},&
 v_i&=s_i+\log s_i,&
 \mu_i&=g_q(\e^{u_i})-g_q(\e^{s_i}),\label{cert:height-data}\\
 d_i&=2s_i^2\mu_i+
  \mathbf{1}_{\{i=M-1\}}\,32c_2q^2\e^{-2s_i}s_i^2. \label{cert:d}
\end{align}
其中 $s_i>0$，$d_i\geq0$。最后一行的附加项
恰来自式 \eqref{dens:B} 中 $16$ 倍的完整余项，
并已乘上正、负高度合并产生的因子 $2$。
遗漏此项会得到不完整的证书。
再定义
\begin{equation}\label{cert:cDE}
 c_j=B_{1,r(j)}W_j(A),\qquad
 D_A=\sum_jB_{2,r(j)}W_j(A),\qquad
 E_A=\sum_jB_{3,r(j)}W_j(A).
\end{equation}
这里 $D_A,E_A$ 只是密度辅助系数，与
第 \ref{sec:trapezoid} 节的差分算子和平凡零点差分不是同一对象。

因为 $s>0$ 时
\[
 \e^{ps}s^{p+2}=s^2\e^{p(s+\log s)},
\]
式 \eqref{dens:Phi} 可改写为
\[
 \frac{\Phi_A(s)}{s^2}
  =\sum_jc_j\e^{p_j(s+\log s)}
          +D_A+\frac{E_A}{s}+\frac{\e^{-A/2+s}}{s}.
\]
由 $d_i$ 的定义逐项展开，可得精确恒等式
\begin{equation}\label{cert:decomposition}
 2\mathcal B(A,q;\boldsymbol u)=\mathcal M+\mathcal A,
\end{equation}
其中
\begin{align}
 \mathcal M&=\sum_{j=0}^{J-1}\sum_{i=0}^{M-1}
                   c_jd_i\e^{p_jv_i},\label{cert:main}\\
 \mathcal A&=\sum_{i=0}^{M-1}d_i
       \left(D_A+\frac{E_A}{s_i}+\frac{\e^{-A/2+s_i}}{s_i}\right).
                                                     \label{cert:aux}
\end{align}
此处没有数值近似，也没有放弃任何一项。
对应
\leanref{M2PilotBudgetIdentity}{weightedPhi_eq_twice_heightBudget}、
\leanref{M2PilotBudgetIdentity}{weightedPhi_eq_main_add_aux}
以及远段的同名恒等式。

\subsection{有余项的 Taylor 不等式}

\begin{lemma}\label{cert:taylor}
对 $0\leq z\leq11$，
\begin{equation}\label{cert:taylor-eq}
 \e^z\leq\sum_{\ell=0}^{20}\frac{z^\ell}{\ell!}
                         +2\frac{z^{21}}{21!}.
\end{equation}
\end{lemma}
\begin{proof}
指数级数中从第 $21$ 项开始，相邻项之比为
$z/(\ell+1)\leq11/22=1/2$。
因此这部分尾和不超过首项 $z^{21}/21!$ 的两倍。
前 $21$ 项保留不变，即得到结论；$z=0$ 也包含在内。
\end{proof}
对应 \leanref{LowRankTaylor}{exp_le_taylor20}。
这不是无余项截断多项式；系数 $2$ 专门吸收全部 Taylor 尾和。

选取移位常数 $p_0,v_0$，满足 $p_j\geq p_0$、$v_i\geq v_0$。
定义四组非负数
\begin{equation}\label{cert:shifted-factors}
 f_j=c_j\e^{p_jv_0},\quad a_j=p_j-p_0,\qquad
 g_i=d_i\e^{p_0(v_i-v_0)},\quad b_i=v_i-v_0 .
\end{equation}
恒等式
\[
 p_jv_i=p_jv_0+p_0(v_i-v_0)+(p_j-p_0)(v_i-v_0)
\]
给出
\begin{equation}\label{cert:factorization}
 \mathcal M=\sum_{j,i}f_jg_i\e^{a_jb_i}.
\end{equation}
若还满足 $a_jb_i\leq11$，则可应用引理 \ref{cert:taylor}。
对非负数列 $(f,a)$ 定义
\[
 M_\ell(f,a)=\sum_j f_ja_j^\ell .
\]
把有限双重和中的多项式逐项展开并交换求和，得到
\begin{equation}\label{cert:low-rank}
 \mathcal M\leq
 \mathcal L_{20}(f,a;g,b):=
 \sum_{\ell=0}^{20}
       \frac{M_\ell(f,a)M_\ell(g,b)}{\ell!}
       +2\,\frac{M_{21}(f,a)M_{21}(g,b)}{21!}.
\end{equation}
原来 $J M$ 个指数项被化为两组各 $22$ 个矩。
这一降维仅使用分配律和有严格余项的 Taylor 不等式；
没有用低秩数值拟合替换原函数。
对应 \leanref{LowRankTaylor}{shifted_exp_factorization}
和 \leanref{LowRankTaylor}{shifted_exp_sum_le}。

\subsection{逐因子向外包围与精确整数矩}

因为 \eqref{cert:low-rank} 中的系数全部非负，
如果四组因子分别向上放大，则右边不会减小。
具体选择非负整数
$F_j,A_j,G_i,B_i$，并逐项证明
\begin{equation}\label{cert:dyadic-bounds}
 0\leq f_j\leq\frac{F_j}{2^{48}},\quad
 0\leq a_j\leq\frac{A_j}{2^{48}},\qquad
 0\leq g_i\leq\frac{G_i}{2^{48}},\quad
 0\leq b_i\leq\frac{B_i}{2^{48}}.
\end{equation}
于是
\[
 M_\ell(f,a)\leq
 \frac{I_\ell}{2^{48(\ell+1)}},
 \quad I_\ell=\sum_j F_jA_j^\ell,
 \qquad
 M_\ell(g,b)\leq
 \frac{J_\ell}{2^{48(\ell+1)}},
 \quad J_\ell=\sum_i G_iB_i^\ell .
\]
从而得到完全有理的上界
\begin{equation}\label{cert:rational}
 \mathcal R=
 \sum_{\ell=0}^{20}
  \frac{I_\ell J_\ell}{\ell!\,2^{96(\ell+1)}}
  +\frac{2I_{21}J_{21}}{21!\,2^{96\cdot22}},
 \qquad \mathcal M\leq\mathcal R.
\end{equation}
欲证明 $\mathcal R\leq r$，只须乘上一个公共正整数分母，
验证一个整数不等式。有限有理运算不会产生浮点舍入误差。

必须区分三个可分别审计的步骤。首先，实数表达式
$f_j,a_j,g_i,b_i$ 由 \eqref{cert:height-data}--\eqref{cert:shifted-factors}
精确定义。其次，\eqref{cert:dyadic-bounds} 是需要证明的不等式，
不能因外部程序输出了某个数就认定它成立。最后，
\eqref{cert:rational} 对已经固定的整数表做精确算术。
本实现的生成脚本负责提出表和包围关系；Lean 的区间引理证明第二步，
内核归约验证第三步。对应
\leanref{LowRankTaylor}{lowRankTaylor20_mono}、
\leanref{LowRankInteger}{lowRankTaylor20_eq_rational}。

具体的有向包围遵守如下规则。例如，若 $s>0$ 且
$\e^\ell\leq s\leq\e^u$ 已获证明，则 $\ell\leq\log s\leq u$。
若 $q\e^{-s}$ 已被一个非负有理区间包围，
利用 $z\mapsto\min\{1,c_1z,c_2z^2\}$ 的单调性，
即可包围 $g_q(\e^s)$。对差
$g_q(\e^{u_i})-g_q(\e^{u_{i+1}})$ 取上界时，
必须用前项的上端减后项的下端；
两个上端相减不构成合法上界。
生成的高度包围证明显式保留前后两个核区间，保证了这一方向。
对 $\exp$ 的有理区间由 Mathlib 的
\texttt{dyadic\_interval} 产生证明项，
而不是接受外部 Arb 计算的成功标记。

完整整数行表位于
\texttt{M2PilotIntegerData} 和 \texttt{M3FarIntegerData}；
逐行包围证明位于相应的
\texttt{Enclosures}、\texttt{FactorBounds} 模块。
例如，中段高度第 $0$ 行最终使用
\[
 G_0=6842520084211003,\qquad B_0=98272331280;
\]
每一行都按同一分母 $2^{48}$ 解释。
这一例子仅说明证书格式，最终不等式使用全部行。
正文给出了所有行的生成对象、核对条件和最终有限算式，
大型整数表则以可逐项读取的源码数据提供。

\subsection{中段试点：区间 [715/4,179]}

中段试点的参数完全固定为
\begin{equation}\label{cert:middle-parameters}
 \begin{gathered}
 A=\frac{715}{4},\quad C=179,\quad k=6,\quad
 q=1+\frac{\e^{179/6}}6,\quad J=1536,\quad M=988,\\
 \sigma_j=\frac{2048+j}{4096}\quad(0\leq j\leq1536),
 \qquad r(j)=\left\lfloor\frac j{128}\right\rfloor,\\
 u_i=\frac{1841+i}{64}\quad(0\leq i\leq988),\qquad
 U=\frac{2829}{64},\qquad
 p_0=\frac13,\quad v_0=\frac{2057}{64}.
 \end{gathered}
\end{equation}
每个左端点的行号为 $0,\ldots,11$，网格在
$1/2$ 开始并在 $7/8$ 结束。高度网格步长为 $1/64$；
在证书双重和中使用右端点 $s_i=u_{i+1}$。
因此不存在把 $988$ 个区间误解为 $988$ 个含首尾端点的节点。

为验证高度条件及 Taylor 区间，只需以下有限标量界：
\begin{gather}
 H\leq\e^{1841/64},\qquad
 q\e^{-1841/64}\geq\frac{31}{64},\label{cert:middle-start}\\
 \e^{215/64}\leq\frac{1842}{64},\qquad
 \frac{2829}{64}\leq\e^{284/64}. \label{cert:middle-log}
\end{gather}
第一个条件保证所有 CS 使用高度合法。第二个条件保证
$c_1q\e^{-u_0}\geq1$，且
$c_2q^2\e^{-2u_0}\geq31>1$，故 $g_q(\e^{u_0})=1$。
式 \eqref{cert:middle-log} 给出
$215/64\leq\log s_i\leq284/64$，于是
\[
 0\leq a_j\leq\frac23,\qquad
 0\leq b_i=s_i+\log s_i-\frac{2057}{64}\leq\frac{33}{2}.
\]
特别地 $a_jb_i\leq11$。这些界分别在
\leanref{M2PilotConstants}{start_logKernel_eq_one}
和 \leanref{M2PilotGrid}{sigma_height_product_bounds} 中落实。

对全部 $1536+988$ 行验证 \eqref{cert:dyadic-bounds} 后，
整数矩算式 \eqref{cert:rational} 满足
\begin{equation}\label{cert:middle-main}
 \mathcal M\leq\mathcal R\leq\frac{3759}{4000}=0.93975.
\end{equation}
实数因子与整数表的连接由
\leanref{M2PilotFactorBounds}{sigmaDyadicBounds}、
\leanref{M2PilotFactorBounds}{heightDyadicBounds} 给出；
有限有理不等式为
\leanref{M2PilotIntegerCertificate}{rationalLowRank20_le}。
该不等式使用普通内核归约
\texttt{decide +kernel}，并非将一个外部数值测试布尔值假定为真。

\subsection{远段试点：区间 [400,1601/4]}

远段使用同一个梯形权和同一个预算结构，但采用四次幂短区间长度：
\begin{equation}\label{cert:far-parameters}
 \begin{gathered}
 A=400,\quad C=\frac{1601}{4},\quad k=4,\quad
 q=1+\frac{\e^{(1601/4)/4}}4,\quad J=384,\quad M=125,\\
 \sigma_j=\frac{512+j}{1024}\quad(0\leq j\leq384),
 \qquad r(j)=\left\lfloor\frac j{32}\right\rfloor,\\
 u_i=\frac{795+i}{8}\quad(0\leq i\leq125),\qquad
 U=115,\qquad p_0=\frac13,\quad v_0=104.
 \end{gathered}
\end{equation}
所需标量界为
\[
 H\leq\e^{795/8},\qquad q\e^{-795/8}\geq\frac{31}{64},
 \qquad \e^{9/2}\leq\frac{796}{8},\qquad115\leq\e^5.
\]
由此 $0\leq a_j\leq2/3$、$0\leq b_i\leq16$，
所以 $a_jb_i\leq32/3<11$。
在相同的 $48$ 位有理格式下，完整的 $384+125$ 行证书给出
\begin{equation}\label{cert:far-main}
 \mathcal M\leq\mathcal R\leq\frac9{20}=0.45.
\end{equation}
对应 \leanref{M3FarGrid}{sigma_height_product_bounds}
和 \leanref{M3FarIntegerCertificate}{rationalLowRank20_le}。
远段证书没有使用历史稿的多项式权；当前远段的全部主张均以
本节的梯形核和式 \eqref{cert:far-parameters} 为准。

\subsection{辅助项、素数幂和平凡零点均已计入}

两个试点中，$B_{2,r}\leq10$、$B_{3,r}\leq168$。
实部权望远镜相消给出
\[
 \sum_jW_j(A)=\e^{-A/8}-\e^{-A/2}\leq\e^{-A/8},
 \quad D_A\leq10\e^{-A/8},\quad E_A\leq168\e^{-A/8}.
\]
高度质量也望远镜相消：
$\sum_i\mu_i=1-g_q(\e^U)\leq1$。
在两个试点中，最后一项
$32c_2q^2\e^{-2U}U^2$ 都不超过 $1$。
所以中段 $\sum_i d_i\leq2(89/2)^2+1<4000$，
远段 $\sum_i d_i\leq2\cdot115^2+1<30000$。
又因 $s_i\geq28$，有
\[
 D_A+\frac{E_A}{s_i}+\frac{\e^{-A/2+s_i}}{s_i}
 \leq16\e^{-A/8}+\frac{\e^{-A/2+U}}{28}.
\]
于是两个辅助预算分别不超过
\begin{align}
 \mathcal A_{\text{中}}
 &\leq4000\left(16\e^{-715/32}+\frac{\e^{-2891/64}}{28}\right)
       \leq\frac1{50000},\label{cert:middle-aux}\\
 \mathcal A_{\text{远}}
 &\leq30000\left(16\e^{-50}+\frac{\e^{-85}}{28}\right)
       \leq\frac1{50000}. \label{cert:far-aux}
\end{align}
最后两个都是有理端点上的指数不等式，另有区间证明，
不是从主项证书的剩余小数位中猜出的误差。
对应 \leanref{M2PilotAuxiliary}{auxiliaryBudget_le_small}
和 \leanref{M3FarAuxiliary}{auxiliaryBudget_le_small}。

由 \eqref{tr:Pcell}，在两个试点各自的 $[A,C]$ 上均有
\begin{equation}\label{cert:smallPV}
 P_k(A,C)\leq\frac1{10^6},\qquad V_A\leq\frac1{10^6}.
\end{equation}
这些估计留有很大余量。例如中段可用
$\epsilon_6(A)\leq1$、$\sqrt{1+\epsilon_6(A)}\leq2$、
$\log(1+\epsilon_6(A))\leq1$、$\log2\geq1/2$，
以及 $\e^{-A/3}\leq\e^{-50}\leq10^{-11}$。
远段的 $\e^{-A/4}=\e^{-100}$ 更小。
平凡零点则直接用 $V_A=\e^{-3A}/(1-\e^{-2A})$。
证明中并未将这些项设成零。

\begin{proposition}[两个完整试点证书]\label{cert:pilots}
在 QRH$(7/8)$、CS v2、Rosser 零点计数和积分显式公式四个输入下，
有以下结论：
\begin{align*}
 \frac{715}{4}\leq\log x\leq179
 &\Longrightarrow \exists p\text{ 为素数},\
                 x<p<x+6x^{5/6},\\
 400\leq\log x\leq\frac{1601}{4}
 &\Longrightarrow \exists p\text{ 为素数},\
                 x<p<x+4x^{3/4}.
\end{align*}
这里 $x>0$，所有区间端点均包含在适用域内。
\end{proposition}
\begin{proof}
中段的完整归一化预算为
\[
 Z+P+V\leq\frac{3759}{4000}+\frac1{50000}
                          +\frac2{10^6}
 =\frac{234943}{250000}=0.939772<\frac{47}{50}<1.
\]
远段的完整预算为
\[
 Z+P+V\leq\frac9{20}+\frac1{50000}+\frac2{10^6}
 =\frac{225011}{500000}=0.450022<\frac12<1.
\]
分别使用命题 \ref{dens:zero-budget} 和命题
\ref{tr:prime-criterion} 即得。
\end{proof}
对应 \leanref{M2PilotResult}{pilot_short_interval_from_inputs}
和 \leanref{M3FarResult}{pilot_short_interval_from_inputs}。
有限高度 RH 尚未用于这两个试点；$H$ 在这里仅是 CS 密度估计的
适用高度。下一节才在向下覆盖中段时使用有限 RH。

\section{连续平移覆盖与有限 RH 下的中段比较}\label{sec:coverage}

两个试点不是仅在有限网格采样点成立的数值观察。
本节证明它们如何覆盖两个完整的实数区间。
向上覆盖通过密度预算的连续平移不等式完成；
向下覆盖 $167\leq\log x\leq715/4$ 时，另使用有限高度 RH
和一个复用原有因子包围的重加权证书。
这样，每一个实数 $\log x$ 都由一条明确的不等式链覆盖，
而不是靠网格点之间的插值假设。

\subsection{实部权和高度项的联合平移}

固定实部网格，并假定
\begin{equation}\label{cover:translation-domain}
 A\geq167,\quad u\geq28,\quad k\geq4,\quad d\geq0,\quad
 \frac12\leq\sigma_j\leq\sigma_{j+1}\leq\frac78,\quad
 \sigma_{j+1}-\sigma_j\leq\frac1{1024}.
\end{equation}
先研究把左端点从 $A$ 改为 $A+d$，
同时把对数高度从 $u$ 改为 $u+d/k$ 的作用。

\begin{lemma}[权的平移界]\label{cover:weight-shift}
在上述条件下，
\begin{equation}\label{cover:weight-shift-eq}
 W_j(A+d)\leq
  \e^{(1/A+\sigma_{j+1}-1)d}W_j(A).
\end{equation}
\end{lemma}
\begin{proof}
由定义，
\[
 W_j(A)=\int_{\sigma_j}^{\sigma_{j+1}}
                         A\e^{(\beta-1)A}\,\dd\beta .
\]
对积分中的每个 $\beta$，利用
$(A+d)/A=1+d/A\leq\e^{d/A}$ 和
$\e^{(\beta-1)d}\leq\e^{(\sigma_{j+1}-1)d}$，
可逐点比较两个非负被积函数。积分后即得结论。
\end{proof}

另一方面，对 $p,r\geq0$，
\begin{equation}\label{cover:height-rate}
 \e^{p(u+d/k)}(u+d/k)^r
 \leq
 \e^{(p+r/u)d/k}\e^{pu}u^r.
\end{equation}
这由 $\log(1+d/(ku))\leq d/(ku)$ 得到。
与 \eqref{cover:weight-shift-eq} 相乘，
一个密度项的总增长指数不超过
\[
 \frac1A+\sigma_{j+1}-1+\frac{p+r/u}{k}.
\]
对于 CS 主项 $p=p_j$、$r=p_j+2$，有
$p_j\leq(8/3)(1-\sigma_j)$、$r\leq3$。
辅助项 $(p,r)=(0,2),(0,1)$ 也满足同样的宽界。
因此在 \eqref{cover:translation-domain} 下，
\begin{align}
 \frac1A+\sigma_{j+1}-1+\frac{p+r/u}{k}
 &\leq\frac1{167}
       -\frac{1-\sigma_{j+1}}3
       +\frac23(\sigma_{j+1}-\sigma_j)+\frac3{112}\notag\\
 &\leq\frac1{167}-\frac1{24}+\frac1{1536}+\frac3{112}<0.
                                                     \label{cover:rate-negative}
\end{align}
这里的严格负上界是一个有限有理不等式；
它解释了为什么限制网格宽度和最低高度后，平移能够覆盖连续参数。

基线项也不会增长。其比值满足
\[
 \frac{\e^{-(A+d)/2+(u+d/k)}(u+d/k)}
      {\e^{-A/2+u}u}
 \leq
 \e^{[-1/2+(1+1/u)/k]d}\leq1.
\]
逐项相加即得以下结论。

\begin{lemma}[密度预算的联合平移]\label{cover:Phi-shift}
在 \eqref{cover:translation-domain} 下，
\begin{equation}\label{cover:Phi-shift-eq}
 \Phi_{A+d}(u+d/k)\leq\Phi_A(u).
\end{equation}
\end{lemma}
对应
\leanref{DensityGridTranslation}{sigmaGridWeight_translation_le}、
\leanref{DensityGridTranslation}{grid_translation_rate_nonpos}
和 \leanref{DensityGridTranslation}{densityGridPhi_translation_le}。
上述证明对实数 $d\geq0$ 直接成立，没有要求 $d$
为某个网格步长的整数倍。

\subsection{完整高度预算的平移不变部分}

令
\begin{equation}\label{cover:shift-parameters}
 q'=q\e^{d/k},\qquad u_i'=u_i+d/k,\qquad U'=U+d/k.
\end{equation}
因为 $q'\e^{-u_i'}=q\e^{-u_i}$，
\begin{equation}\label{cover:kernel-invariant}
 g_{q'}(\e^{u_i'})=g_q(\e^{u_i}),\qquad
 q'^2\e^{-2U'}=q^2\e^{-2U}.
\end{equation}
故高度质量 $\mu_i$ 和余项系数均保持不变。
再用引理 \ref{cover:Phi-shift} 逐项比较
式 \eqref{dens:B}，得到
\begin{proposition}[完整预算向上平移]\label{cover:B-shift}
若实部网格宽度不超过 $1/1024$，$A\geq167$、
所有 $u_i\geq28$、$k\geq4$，且原高度网格合法，则
\begin{equation}\label{cover:B-shift-eq}
 \mathcal B(A+d,q\e^{d/k};\boldsymbol u+d/k)
 \leq\mathcal B(A,q;\boldsymbol u)\qquad(d\geq0).
\end{equation}
平移后的高度网格仍满足
\eqref{dens:height-conditions}。
\end{proposition}
\begin{proof}
高度下界、单调性和 $U\geq6$ 在非负平移后保持；
初始核等于 $1$ 由 \eqref{cover:kernel-invariant} 保持。
式 \eqref{dens:B} 的所有系数非负，
而每个 $\Phi$ 因子由 \eqref{cover:Phi-shift-eq} 不增，
故包含余项的完整和不增。
\end{proof}
对应
\leanref{TrapezoidBudgetTranslation}{trapezoidHeightBudget_translation_le}。

还必须验证 $q'$ 足够大，不能只平移一个与实际区间无关的核。
由 \eqref{tr:domains}，
\begin{equation}\label{cover:q-valid}
 q_k(C+d)=1+\frac{\e^{(C+d)/k}}k
 \leq\e^{d/k}\left(1+\frac{\e^{C/k}}k\right)=q_k(C)\e^{d/k}.
\end{equation}
若目标点 $L=A+d$ 且试点右端点 $C\geq A$，
则 $q_k(L)\leq q_k(C+d)\leq q'$。
所以命题 \ref{dens:zero-budget} 对真实长度 $h_k(L)$ 适用。
我们只要求核参数是一个上界，并不要求恰好等于 $q_k(L)$。

\subsection{上中段与远段的连续覆盖}

对第 \ref{sec:certificates} 节的中段试点，
令 $A_0=715/4$。给定
$A_0\leq L\leq400$，取 $d=L-A_0$，
将试点高度网格平移 $d/6$，并令 $q'=q\e^{d/6}$。
由命题 \ref{cover:B-shift} 和式 \eqref{cert:middle-main}、
\eqref{cert:middle-aux}，
\begin{equation}\label{cover:upper-zero}
 \frac{|Z_T(\e^L,h_6(L))|}{h_6(L)}
 \leq2\mathcal B(A_0,q;\boldsymbol u)
 \leq\frac{93977}{100000}.
\end{equation}
此式对所有 $T$、所有该范围内的实数 $L$ 都成立。
取单元端点 $A=L$、$C=L+1/4$，
有 $C\leq401$。在这一范围，
\begin{equation}\label{cover:upper-small}
 P_6(L,L+1/4)\leq10^{-6},\qquad V_L\leq10^{-6}.
\end{equation}
例如用 $\epsilon_6(L)\leq1$、
$\sqrt{1+\epsilon_6(L)}\leq2$、
$\log(1+\epsilon_6(L))\leq1$ 和 $\log2\geq1/2$，
式 \eqref{tr:Pcell} 的指数前因子不超过 $200000$。
又有 $\e^{-L/3}\leq\e^{-55}\leq10^{-18}$，
已足以证明第一式。第二式由平凡零点预算的单调性得到。
所以 $Z+P+V\leq0.939772<1$，
覆盖整个 $715/4\leq L\leq400$。

对远段试点，令 $A_0=400$。给定 $400\leq L\leq700$，
取 $d=L-400$，高度平移 $d/4$，
其余步骤相同，得到
\[
 \frac{|Z_T(\e^L,h_4(L))|}{h_4(L)}
 \leq\frac{22501}{50000}.
\]
这次 $C=L+1/4\leq701$，并且
$P_4(L,L+1/4),V_L\leq10^{-6}$。
例如指数前因子可放大到 $500000$，而
$\e^{-L/4}\leq\e^{-100}\leq10^{-18}$。
因此完整预算不超过 $0.450022<1$。

\begin{proposition}[两个向上覆盖区间]\label{cover:upper-ranges}
在 QRH$(7/8)$、CS v2、Rosser 零点计数和积分显式公式下，对每个 $x>0$，
\begin{align*}
 \frac{715}{4}\leq\log x\leq400
 &\Longrightarrow\exists p\text{ 为素数},\
                         x<p<x+6x^{5/6},\\
 400\leq\log x\leq700
 &\Longrightarrow\exists p\text{ 为素数},\
                         x<p<x+4x^{3/4}.
\end{align*}
\end{proposition}
由前面的预算和命题 \ref{tr:prime-criterion} 即得。
对应
\leanref{M3MiddleUpper}{middle_upper_short_interval_from_inputs}
和 \leanref{M3FarUpper}{far_short_interval_from_inputs}。
后一个结论是四次幂长度在一个有限的对数范围内成立，
不应改写为四次勒让德问题的全范围结论。

\subsection{有限 RH 怎样改善下中段的核}

向下平移时 \eqref{cover:Phi-shift-eq} 的方向不能直接反用。
我们改用有限 RH 输入 \cite{PlattTrudgian}：
$0<|\gamma|\leq H$ 的非平凡零点均有 $\beta=1/2$。
设 $0\leq r\leq1$、$q\leq rQ$，并且
\begin{equation}\label{cover:linear-threshold}
 \frac{c_1q}{H}\leq r.
\end{equation}
则对于 $t\geq H$，
\begin{equation}\label{cover:scaled-kernel}
 g_q(t)\leq r g_Q(t).
\end{equation}
确实，$g_q(t)\leq c_1q/t\leq r$；
又分别有 $c_1q/t\leq r c_1Q/t$ 和
$c_2q^2/t^2\leq r^2c_2Q^2/t^2\leq r c_2Q^2/t^2$。
这三个上界一起给出 \eqref{cover:scaled-kernel}。
仅有 $q\leq rQ$ 并不足以缩放核的常数分支 $1$；
条件 \eqref{cover:linear-threshold} 是必要的额外检查。

把正高度零点分成 $0<\gamma\leq H$ 与 $\gamma>H$。
第一部分由有限 RH 和 $g_q\leq1$ 控制；
第二部分用 \eqref{cover:scaled-kernel}，再扩大为全部正高度的非负和。
由 Rosser 上界 $N(H)\leq H\log H$ 及
式 \eqref{dens:positive-final}，
\begin{equation}\label{cover:low-high}
 \sum_{0<\gamma\leq T}m_\rho\e^{(\beta-1)L}g_q(\gamma)
 \leq\e^{-L/2}H\log H
             +r\mathcal B(L,Q;\boldsymbol u).
\end{equation}
因此完整零点预算可取
\begin{equation}\label{cover:lower-Z}
 Z_{\text{低中}}(L)=
 2\e^{-L/2}H\log H+2r\mathcal B(L,Q;\boldsymbol u).
\end{equation}
有限 RH 只用于第一项；高于 $H$ 的零点仍使用 QRH 和 CS 密度估计。
对应
\leanref{TrapezoidLowHeight}{norm_truncatedTrapezoidZeroSum_le_scaled_heightBudget}。

\subsection{缩放参数的合法性}

现在限制
$167\leq L\leq A_0=715/4$，保留中段试点的固定 $Q=q$ 和高度网格，
并取
\[
 r=\e^{(L-A_0)/6},\qquad q=q_6(L).
\]
显然 $0<r\leq1$。为避免同名混淆，本小节的 $Q$ 始终是
试点常数 $1+\e^{179/6}/6$，$q$ 则随 $L$ 变化。
由直接恒等变形，
\[
 \frac{q_6(L)}r
   =\e^{(A_0-L)/6}+\frac{\e^{A_0/6}}6
   \leq \e^{47/24}+\frac{\e^{715/24}}6 .
\]
以下两个标量不等式已用有理区间核对：
\begin{equation}\label{cover:cap}
 \e^{47/24}+\frac{\e^{715/24}}6
 \leq Q,\qquad
 \e^{47/24}+\frac{\e^{715/24}}6
 \leq1453125000000=\frac{H}{c_1}.
\end{equation}
因此 $q_6(L)\leq rQ$ 且 $c_1q_6(L)/H\leq r$，
正好满足 \eqref{cover:linear-threshold}。
这些并非渐近比较，而是包含下端点 $L=167$ 的统一界。
对应 \leanref{M3LowerConstants}{middleQ_le_scaling_q}
和 \leanref{M3LowerConstants}{middleQ_linear_le_scaling}。

\subsection{三档重加权，而非把全主项乘以粗常数}

还需控制 \eqref{cover:lower-Z} 中的
$2r\mathcal B(L,Q;\boldsymbol u)$。
直接将试点主项整体乘以 $64$ 会失去结论，
所以必须保留各个实部单元的差异。

记 $d=A_0-L\in[0,47/4]$。由实部权的积分表示，
\begin{equation}\label{cover:weight-down}
 rW_j(L)\leq
       \e^{(5/6-\sigma_j)d}W_j(A_0).
\end{equation}
证明如下：将 $rL\e^{(\beta-1)L}$ 与
$A_0\e^{(\beta-1)A_0}$ 相比，前者与后者的比为
$(L/A_0)\e^{(5/6-\beta)d}$。因 $L/A_0\leq1$、
$\beta\geq\sigma_j$，该比不超过
$\e^{(5/6-\sigma_j)d}$；积分即得。

在中段的 $1536$ 个实部单元上，定义
\begin{equation}\label{cover:multiplier}
 \lambda_j=
 \begin{cases}
 64,&0\leq j<1229,\\
 3/2,&1229\leq j<1366,\\
 1,&1366\leq j<1536.
 \end{cases}
\end{equation}
这些分界由左端点决定：
$j\geq1229$ 蕴含 $\sigma_j\geq4/5$，
$j\geq1366$ 蕴含 $\sigma_j\geq5/6$。
在第一档，$5/6-\sigma_j\leq1/3$，
所以指数因子不超过 $\e^{47/12}\leq64$；
在第二档不超过 $\e^{47/120}\leq3/2$；
在第三档不超过 $1$。
故
\begin{equation}\label{cover:weighted-bound}
 rW_j(L)\leq\lambda_jW_j(A_0).
\end{equation}
越过 $4/5$ 或 $5/6$ 的单元仍按它的左端点处理，
不会把跨界单元错误地分到较小系数一侧。
基线满足
\[
 r\e^{-L/2+u}u
  =\e^{(A_0-L)/3}\e^{-A_0/2+u}u
  \leq64\e^{-A_0/2+u}u.
\]
CS 的两个辅助部分也可统一乘以 $64$，
因为其原试点预算已很小。

在式 \eqref{cert:main} 中只把实部系数 $c_j$
替换为 $\lambda_jc_j$，记所得主项为
\begin{equation}\label{cover:weighted-main}
 \mathcal M_{\lambda}
   =\sum_{j,i}\lambda_jc_jd_i\e^{p_jv_i}.
\end{equation}
这里 $c_j,d_i,p_j,v_i$ 全部保持中段试点在 $A_0$ 的原值。
由精确拆分 \eqref{cert:decomposition}、非负性及
\eqref{cover:weighted-bound}，可逐项得到
\begin{equation}\label{cover:scaled-B}
 2r\mathcal B(L,Q;\boldsymbol u)
     \leq\mathcal M_\lambda+64\mathcal A_{\text{中}}.
\end{equation}
这是下中段重加权的关键表达式。
对应 \leanref{M3LowerComparison}{scaled_sigma_weight_le}
和 \leanref{M3LowerExpression}{scaled_heightBudget_le}。

\subsection{复用整数表的第三个有限算术检查}

式 \eqref{cover:weighted-main} 使用的 $a_j,b_i,g_i$
及原有实数包围全部不变，仅将 $f_j$ 乘以 $\lambda_j$。
因此第 \ref{sec:certificates} 节的 Taylor 矩论证仍然适用。
令
\[
 m_j=2\lambda_j=
 \begin{cases}
 128,&j<1229,\\
 3,&1229\leq j<1366,\\
 2,&j\geq1366.
 \end{cases}
\]
在中段整数表中，把每个实部行 $(F_j,A_j)$ 改为
$(m_jF_j,A_j)$；高度行 $(G_i,B_i)$ 完全不动。
这仍是一个纯整数表，且其低秩算式是
$2\mathcal M_\lambda$ 的上界。精确有理检查给出
\begin{equation}\label{cover:weighted-certificate}
 \mathcal R_{\text{加倍重权}}\leq\frac{193}{100},
 \qquad
 \mathcal M_\lambda\leq\frac{193}{200}=0.965.
\end{equation}
第二式使用低秩矩算式对第一组系数的一次齐次性。
不需要重新信任一个浮点运行，也没有新增外部解析输入。
对应 \leanref{M3LowerIntegerCertificate}{rational_double_budget}、
\leanref{M3LowerInteger}{weighted_lowRank_le}
和 \leanref{M3LowerMainBound}{weightedMain_le}。

对 $L\geq167$，有限 RH 低零点预算满足
\begin{equation}\label{cover:low-small}
 2\e^{-L/2}H\log H\leq10^{-6}.
\end{equation}
例如 $\e^{-L/2}\leq\e^{-80}$、$\log H\leq H$，
而 $2H^2\e^{-80}\leq10^{-6}$；这是一个很宽的有理区间检查。
在 $167\leq L\leq715/4$ 上，
$P_6(L,L)\leq10^{-6}$ 和 $V_L\leq10^{-6}$，
可由第 \ref{sec:trapezoid} 节的公式直接推出。
结合 \eqref{cover:scaled-B}、\eqref{cover:weighted-certificate}
和 $\mathcal A_{\text{中}}\leq1/50000$，完整预算为
\begin{equation}\label{cover:lower-complete}
 Z_{\text{低中}}(L)+P_6(L,L)+V_L
 \leq\frac{193}{200}+\frac{64}{50000}+\frac3{10^6}
 =0.966283<1.
\end{equation}
三份 $10^{-6}$ 分别是低高度零点、真素数幂和平凡零点，
并非重复计入同一个误差项。

\begin{proposition}[下中段覆盖]\label{cover:lower-range}
在 QRH$(7/8)$、CS v2、Rosser 零点计数、积分显式公式和
有限 RH 到 $H=3\cdot10^{12}$ 五个输入下，对每个 $x>0$，
\[
 167\leq\log x\leq\frac{715}{4}
 \quad\Longrightarrow\quad
 \exists p\text{ 为素数},\quad x<p<x+6x^{5/6}.
\]
\end{proposition}
\begin{proof}
式 \eqref{cover:lower-Z} 是对任意截断高度成立的零点预算；
式 \eqref{cover:lower-complete} 验证三项判据。
用命题 \ref{tr:prime-criterion} 即得。
\end{proof}
对应 \leanref{M3LowerResult}{lower_middle_short_interval_from_inputs}。

\subsection{本节的适用范围及其与全范围结论的衔接}

合并命题 \ref{cover:upper-ranges} 与
\ref{cover:lower-range}，当前证明得到
\[
 \begin{array}{c|c|c}
 \log x\text{ 的闭区间}&\text{保证含素数的开区间}
                                  &\text{是否使用有限 RH}\\ \hline
 {[167,715/4]}&(x,x+6x^{5/6})&\text{是}\\
 {[715/4,400]}&(x,x+6x^{5/6})&\text{否}\\
 {[400,700]}&(x,x+4x^{3/4})&\text{否}
 \end{array}
\]
公共端点在相邻论证中均被包含，所以不存在接缝空隙。
远段中 $x>1$，故 $4x^{3/4}<6x^{5/6}$，
其结论也足以提供六次幂短区间所需的素数。
小于 $\e^{167}$ 的范围及大于等于 $\e^{700}$ 的无穷尾段
由后续章节的独立分支处理；本节不把有限区间证明扩张为全范围
$k=4$ 结论。

历史上逐单元认证的记录仍有复现和研究价值，但当前数学覆盖的核心是
本节的连续参数不等式以及两个试点、一次重加权的有限算术检查。
因此，早期证书的敏感度数字不能自动作为当前证书的统一余量。
当前可以直接从正文核对的三段上界分别为
$0.966283$、$0.939772$ 和 $0.450022$；
它们仅描述本节这些固定输入和固定证书下的预算。

\section{无限尾部：\texorpdfstring{$L\geq700$}{L>=700}}
\label{sec:tail}

本节证明一个适用于所有充分大实数的严格短区间结论。它只使用
H1（QRH）、H2（CS v2）、H3（Rosser 零点计数）和 H6（CHJII
短区间显式公式）。H4 的积分显式公式、H5 的有限高度 RH 以及 H7
的 CHL 行均不参与本节的推导。特别地，以下出现
$3\cdot10^{12}$ 是为了满足 CS 密度估计的适用域，并非在尾部另用一次有限 RH。

与有限解析段不同，本节直接估计未平滑的 $\psi$ 增量。需要完成的连接是：
CHJII 的同一截断高度、实际零点和的复数模、密度积分、素数幂扣除，
以及将右闭区间改为严格开区间。六项预算分别承担这些任务；其中任何一项
都不是以一个未证的局部素数结论代替。

\subsection{参数、定义域与同一截断高度}

令
\begin{equation}\label{tail:parameters}
 x=\e^L,\qquad h=4\e^{3L/4},\qquad y=x+h,\qquad
 T=\e^{9L/32},\qquad V=2T,
\end{equation}
并记
\begin{equation}\label{tail:u-lambda}
 u=\log V=\frac9{32}L+\log2,\qquad
 \lambda=L-\frac83\log V=\frac14L-\frac83\log2.
\end{equation}
这里 $V$ 是一个确定的高度上界；CHJII 给出的实际截断高度将记为 $T^*$。
二者不可混淆。以下所有估计均在 $L\geq700$ 上成立。

\begin{lemma}[尾部参数的合法性]\label{tail:domain}
上述参数满足
\begin{align}
 &\max\{51,L^2\}<T<\frac{\sqrt{x}-2}{4},\label{tail:T-domain}\\
 &V\geq3\cdot10^{12},\qquad u>1,\qquad \lambda>0,\label{tail:density-domain}\\
 &0<h\leq x,\qquad x<y\leq2x,\qquad
 L\leq\log y\leq L+1,\qquad \sqrt y\leq2\sqrt x.
 \label{tail:elementary-domain}
\end{align}
\end{lemma}

\begin{proof}
由 $\e^t\geq1+t$，有 $T>51$。三阶指数下界给出
\[
 T\geq\frac{(9L/32)^3}{3!}
 =\frac{729L^3}{196608}>L^2,
\]
最后一步用 $L\geq700$。另一方面，
\[
 \e^{L/2}=T\e^{7L/32},\qquad \e^{7L/32}>6.
\]
结合 $T>51$ 得 $\e^{L/2}>4T+2$，即式~\eqref{tail:T-domain}。
由于 $9L/32\geq100$，再用
$\e^{100}\geq100^{10}/10!>3\cdot10^{12}$ 即可验证高度下界。
式~\eqref{tail:u-lambda} 与 $0<\log2<1$ 给出 $u>1$ 和 $\lambda>0$。

由 $\e^{L/4}\geq4$ 得 $h\leq x$，因而 $y\leq2x$。
取对数得到 $\log y\leq L+\log2\leq L+1$；下界来自 $y\geq x$。
最后，$y\leq2x\leq4x$ 且平方根非负，故 $\sqrt y\leq2\sqrt x$。
所有这些估计都是整个半直线上的不等式，并非基点处的数值观察。
\end{proof}

对应实现为 \leanref{TailDomains}{tail_log_sq_lt_T}、
\leanref{TailDomains}{tail_T_lt_sqrt_limit}、
\leanref{TailDomains}{tail_height_ge_RH} 和同文件的端点估计。

以 $m_\rho$ 表示实际 $\zeta$ 零点 $\rho=\beta+\ii\gamma$ 的解析重数，定义
\begin{equation}\label{tail:actual-sum}
 S(x,y;W)=\sum_{\substack{\zeta(\rho)=0,\ 0<\beta<1\\|\gamma|\leq W}}
 m_\rho\frac{y^\rho-x^\rho}{\rho}.
\end{equation}
这是有限和，截断条件为闭条件 $|\gamma|\leq W$。正实底数的复幂取
$t^\rho=\exp(\rho\log t)$，所以不存在分支选择问题。

H6 是 \cite{CHJII} 的定理 1.3 与表 1 首行的明确特例：
$\log x_M=40$、$\alpha=1/2$、$\omega=0$、$M=2.894=1447/500$。
由引理~\ref{tail:domain}，存在\emph{同一个}
$T^*\in[T,2T]$，使
\begin{equation}\label{tail:chjii}
 \left|\psi(y)-\psi(x)-h+S(x,y;T^*)\right|
 \leq\frac{1447}{500T}\bigl(y\log y+x\log x\bigr).
\end{equation}
不能把两个单点显式公式各自存在的截断高度直接合并成式~\eqref{tail:chjii}。
这里采用的就是源文的短区间版本；其量词和复数模形式明确保存在
\leanref{LiteratureInputs}{CHJIIShortInterval} 中。

\subsection{从复数零点增量到带重数的正实和}

\begin{lemma}[复幂增量]\label{tail:cpow-increment}
若 $0<x\leq y$、$\rho\ne0$ 且 $\beta=\Re\rho\leq1$，则
\begin{equation}\label{tail:cpow-bound}
 \left|\frac{y^\rho-x^\rho}{\rho}\right|
 \leq (y-x)x^{\beta-1}.
\end{equation}
\end{lemma}

\begin{proof}
函数 $t\mapsto t^\rho/\rho$ 在正实轴上的导数为 $t^{\rho-1}$，
其模为 $t^{\beta-1}$。由于 $\beta-1\leq0$，在 $[x,y]$ 上该模
不超过 $x^{\beta-1}$。对复值函数使用导数范数的区间估计，即得结论。
有限求和后，三角不等式与 $m_\rho\geq0$ 保留不等号方向。
\end{proof}

这一论证见 \leanref{CpowInterval}{norm_cpow_sub_div_le} 和
\leanref{CpowInterval}{norm_sum_cpow_sub_div_le}。于是
\begin{equation}\label{tail:sum-reduction}
 |S(x,y;T^*)|
 \leq h\sum_{|\gamma|\leq T^*}m_\rho\e^{(\beta-1)L}
 \leq h\sum_{|\gamma|\leq V}m_\rho\e^{(\beta-1)L}.
\end{equation}
第二个不等式只使用 $T^*\leq V$ 和各项非负，并没有要求所选的 $T^*$
避开某个零点的高度。

实际零点集合在共轭下不变，共轭零点的解析重数相等。
本项目还证明了实轴 $0<\beta<1$ 上没有零点。因此
\begin{equation}\label{tail:pairing}
 \sum_{|\gamma|\leq V}m_\rho\e^{(\beta-1)L}
 =2\sum_{0<\gamma\leq V}m_\rho\e^{(\beta-1)L}.
\end{equation}
这里的因子 $2$ 依赖这些已经证明的对称性与实轴非零性，
不能仅由“通常零点成对”一句话隐去可能的实轴项。
相应定理为 \leanref{ZetaZeros}{sum_allZetaZeros_eq_twice}。

\subsection{分层积分、Rosser 计数与 CS 包围}

用 $N(V)$ 表示 $0<\gamma\leq V$ 的带重数零点数，
用 $N(\sigma,V)$ 表示其中还满足 $\sigma\leq\beta$ 的零点数。
两处高度条件都是右闭的，实部条件也是闭的。

\begin{lemma}[有限零点分层]\label{tail:layer}
在 H1 下，对 $L\geq0$ 有
\begin{equation}\label{tail:layer-formula}
 \sum_{0<\gamma\leq V}m_\rho\e^{(\beta-1)L}
 \leq N(V)\e^{-3L/8}
 +\int_{5/8}^{7/8}L\e^{(\sigma-1)L}N(\sigma,V)\,\dd\sigma.
\end{equation}
\end{lemma}

\begin{proof}
H1 给出所有相关零点的 $\beta\leq7/8$。
若 $\beta\leq5/8$，则 $\e^{(\beta-1)L}\leq\e^{-3L/8}$。
若 $5/8\leq\beta\leq7/8$，则微积分基本定理给出
\[
 \e^{(\beta-1)L}=\e^{-3L/8}
 +\int_{5/8}^{7/8}
  \mathbf{1}_{\{\sigma\leq\beta\}}L\e^{(\sigma-1)L}\,\dd\sigma.
\]
乘以重数并对有限个零点求和即可。
每个被积函数是连续函数乘以一个可测区间的示性函数，因而区间可积；
这里只交换有限和与积分，不涉及未经说明的无穷交换。
\end{proof}

一般有限和版本见 \leanref{FiniteZeroLayer}{finite_exp_layer_bound}，
实际零点版本见 \leanref{TailDensity}{positive_zero_exp_layer}。

H3 保留了 Rosser 计数式的全部常数及原来的闭矩形计数约定
\cite{Rosser1941,Dusart2016}。若该计数记为 $N_R(V)$，则
\begin{align}
 N(V)&\leq N_R(V),\notag\\
 \left|N_R(V)-\left(\frac{V}{2\pi}\log\frac{V}{2\pi}
       -\frac{V}{2\pi}+\frac78\right)\right|
 &\leq0.137\log V+0.443\log\log V+1.588.
 \label{tail:rosser}
\end{align}
前一个不等式是计数集合的包含关系；我们没有把外部计数集合默默改成另一集合。
当 $V\geq3\cdot10^{12}$ 且 $\log V\geq1$ 时，由
$\log\log V\leq\log V$、$\pi\geq1$ 等初等比较，得到
\begin{equation}\label{tail:coarse-count}
 N(V)\leq V\log V.
\end{equation}
例如，将式~\eqref{tail:rosser} 中负项丢去后，主项至多为
$(V/2)\log V$，其余至多为
$0.580\log V+2.463\leq3.043\log V$；在本域中这小于剩余的
$(V/2)\log V$。因此式~\eqref{tail:coarse-count} 是从 H3 推出的粗界，
不是第八项外部前提。
Lean 对应为 \leanref{LiteratureInputs}{totalZeroCount_le_mul_log}。

在 $5/8\leq\sigma\leq7/8$ 上，H2 的相关行具有
\[
 B_1\leq\frac{561}{25}=22.44,\qquad
 B_2\leq\frac{829}{100}=8.29,\qquad B_3\leq147.
\]
若 $p(\sigma)=3(1-\sigma)/(2-\sigma)$，则
\[
 p(\sigma)\leq\frac83(1-\sigma),\qquad
 \frac{7-5\sigma}{2-\sigma}=p(\sigma)+2\leq3.
\]
结合 $V\geq1$ 和 $\log V\geq1$，可从 CS v2 的原指数直接推出
\begin{equation}\label{tail:cs-envelope}
 N(\sigma,V)\leq
 \frac{561}{25}V^{\frac83(1-\sigma)}u^3
 +\frac{829}{100}u^2+147u.
\end{equation}
这里的 $\log^3 V$ 是本节\emph{由 v2 原式向上放宽后}的统一包围；
它不把原始 CS 输入改写为另一个版本。
所用行的存在性、常数比较及指数比较见
\leanref{DensityInput}{cs_tail_density}。

令 $C=(561/25)u^3$、$D=(829/100)u^2+147u$。
由于
\[
 \e^{(\sigma-1)L}V^{\frac83(1-\sigma)}
 =\e^{(\sigma-1)\lambda},
\]
式~\eqref{tail:cs-envelope} 对应的积分可以精确计算为
\begin{align}
 &\int_{5/8}^{7/8}L\e^{(\sigma-1)L}
 \left(CV^{\frac83(1-\sigma)}+D\right)\,\dd\sigma\notag\\
 &\quad=C\frac{L}{\lambda}
      \left(\e^{-\lambda/8}-\e^{-3\lambda/8}\right)
       +D\left(\e^{-L/8}-\e^{-3L/8}\right)\notag\\
 &\quad\leq C\frac{L}{\lambda}\e^{-\lambda/8}+D\e^{-L/8}.
 \label{tail:density-integral}
\end{align}
最后一步先用 $C,D,L\geq0$、$\lambda>0$ 确认系数非负，再舍去两个
负的下端点项。对应定理为
\leanref{TailDensity}{integral_cs_tail_envelope_le}。

把式~\eqref{tail:coarse-count}、\eqref{tail:density-integral}
代入分层界，并用式~\eqref{tail:pairing} 乘以 $2$，定义
\begin{align}
 Z_0(L)&=4u\e^{-3L/32},\label{tail:Z0}\\
 Z_1(L)&=\frac{1122}{25}u^3\frac{L}{\lambda}\e^{-\lambda/8},
 \label{tail:Z1}\\
 Z_2(L)&=2\left(\frac{829}{100}u^2+147u\right)\e^{-L/8}.
 \label{tail:Z2}
\end{align}
第一项的指数来自
$2V u\e^{-3L/8}=4u\e^{-3L/32}$，其中 $V=2\e^{9L/32}$。
于是对 CHJII 选择的任何合法 $T^*$ 都有
\begin{equation}\label{tail:zero-bound}
 |S(x,y;T^*)|\leq h\bigl(Z_0(L)+Z_1(L)+Z_2(L)\bigr).
\end{equation}
这条从实际零点和到三项实数预算的完整连接是
\leanref{TailZeroSum}{tail_truncated_zero_bound}。

\subsection{截断误差、素数幂与严格右端点}

记后三项预算为
\begin{align}
 R(L)&=\frac{1447}{2000}(3L+2)\e^{-L/32},\label{tail:R}\\
 P(L)&=\frac12(L+1)^2\e^{-L/4},\label{tail:P}\\
 E(L)&=\frac14(L+1)\e^{-3L/4}.\label{tail:E}
\end{align}
由 $y\leq2x$ 及 $\log y\leq L+1$，
\[
 y\log y+x\log x\leq x(3L+2).
\]
除以区间长度 $h$ 后，CHJII 的右端满足
\begin{equation}\label{tail:R-proof}
 \frac{1447}{500T}(y\log y+x\log x)
 \leq\frac{1447}{500T}x(3L+2)=hR(L).
\end{equation}
等号中的指数核算为 $x/T=\e^{23L/32}$ 和
$h\e^{-L/32}=4\e^{23L/32}$。

\begin{lemma}[素数幂与右端点的预算]\label{tail:arithmetic-errors}
在式~\eqref{tail:parameters} 的参数域上，
\begin{equation}\label{tail:arithmetic-estimates}
 \psi(y)-\theta(y)\leq hP(L),\qquad \log y\leq hE(L).
\end{equation}
\end{lemma}

\begin{proof}
Mathlib 已证明的无条件 Chebyshev 界给出
$\psi(y)-\theta(y)\leq2\sqrt y\log y$。
由引理~\ref{tail:domain}，
\[
 2\sqrt y\log y\leq4\sqrt x(L+1)
 \leq2\sqrt x(L+1)^2=hP(L).
\]
这里第二步只需 $L+1\geq2$。另一方面，式~\eqref{tail:E} 给出
$hE(L)=L+1$，所以 $\log y\leq hE(L)$。
\end{proof}

这三条预算分别对应
\leanref{TailErrorBudgets}{tail_truncation_error_le}、
\leanref{TailErrorBudgets}{tail_prime_power_error_le} 和
\leanref{TailErrorBudgets}{tail_endpoint_error_le}。
此处采用实际 $\psi-\theta$ 的已证明上界，没有把“素数幂很少”当作直觉步骤。

若 $0\leq x\leq y$、$y\geq1$ 且开区间 $(x,y)$ 中没有素数，
则 $\theta(y)-\theta(x)$ 的和中至多只有右端点 $y$ 本身这一项，因而
\begin{equation}\label{tail:open-extraction}
 \theta(y)-\theta(x)\leq\log y.
\end{equation}
无论 $y$ 是整数、素数还是非整数，这个上界都成立。
故严格不等式 $\theta(y)-\theta(x)>\log y$ 足以保证开区间内有素数。
形式化中的对应引理为
\leanref{PrimePowerBudget}{prime_of_theta_sub_gt_log}。

\subsection{六项预算的基点证书与全域单调性}

定义完整归一化预算
\begin{equation}\label{tail:delta}
 \mathcal D(L)=Z_0(L)+Z_1(L)+Z_2(L)+R(L)+P(L)+E(L).
\end{equation}
它就是 \leanref{TailArithmetic}{tailBound}，没有省略低高度零点、
负高度零点、素数幂或严格端点费用。

\begin{lemma}[基点的内核数值界]\label{tail:base-certificate}
有 $\mathcal D(700)<14/25$。
\end{lemma}

\begin{proof}
反射式区间证明首先建立
$0.693147<\log2<0.693148$，进而得到
\[
 197.568147<u(700)<197.568148,\qquad
 173.1516<\lambda(700)<173.1517.
\]
为避免直接对正分母作粗糙除法，将 $Z_1$ 之外的五项之和记为 $A$，
并记
\[
 B=44.88\,u(700)^3\,700\,\e^{-\lambda(700)/8}.
\]
于是 $\mathcal D(700)=A+B/\lambda(700)$。
源文件中的精确有理算术和 64 位二进区间反射证明验证
\[
 A\lambda(700)+B<0.56\lambda(700).
\]
因 $\lambda(700)>0$，除去该正因子便得到所述结论。
这些有限小数均被解释为有理数；指数包围的证明项由 Lean 内核检查，
不是读取 Python 或 Arb 的成功标志。
\end{proof}

对应定理为 \leanref{TailArithmetic}{log_two_bounds} 和
\leanref{TailArithmetic}{tail_at_700}。
基点计算本身还不足以证明整个无限尾部；以下单调性补全这一点。

\begin{lemma}[连续单调性]\label{tail:monotone}
$\mathcal D$ 在 $[700,\infty)$ 上单调不增，因此
\begin{equation}\label{tail:global-budget}
 L\geq700\quad\Longrightarrow\quad
 \mathcal D(L)\leq\mathcal D(700)<\frac{14}{25}<1.
\end{equation}
\end{lemma}

\begin{proof}
对 $a,c\geq0$、整数 $m\geq1$，令
$f(L)=(aL+b)^m\e^{-cL+d}$。在 $aL+b\geq0$ 时，
\[
 f'(L)=(aL+b)^{m-1}\e^{-cL+d}
       \bigl(ma-c(aL+b)\bigr).
\]
所以，只要 $ma\leq c(700a+b)$，便有 $f'(L)\leq0$ 对所有 $L\geq700$
成立。将 $b_0=\log2$，$0<b_0<1$，代入下表的因子，即可逐项验证这一条件：
\begin{center}
\begin{tabular}{c|c|c|c}
所属项 & $aL+b$ & $m$ & $c$\\\hline
$Z_0$ & $9L/32+b_0$ & $1$ & $3/32$\\
$Z_1$ 的指数多项式因子 & $9L/32+b_0$ & $3$ & $1/32$\\
$Z_2$ 的二次部分 & $9L/32+b_0$ & $2$ & $1/8$\\
$Z_2$ 的一次部分 & $9L/32+b_0$ & $1$ & $1/8$\\
$R$ & $3L+2$ & $1$ & $1/32$\\
$P$ & $L+1$ & $2$ & $1/4$\\
$E$ & $L+1$ & $1$ & $3/4$
\end{tabular}
\end{center}
第二行用的是
$\e^{-\lambda/8}=\e^{-L/32+b_0/3}$。
它剩下的因子 $L/\lambda$ 也单调不增：
\[
 \frac{\dd}{\dd L}\frac{L}{L/4-(8/3)b_0}
 =-\frac{(8/3)b_0}{(L/4-(8/3)b_0)^2}\leq0.
\]
由于相关因子均非负，乘积与非负线性组合仍单调不增。
这证明了全域单调性，再结合基点界即得式~\eqref{tail:global-budget}。
\end{proof}

该导数判据、比例因子及组合步骤分别在
\leanref{TailMonotonic}{affine_power_exp_antitone}、
\leanref{TailMonotonic}{tail_ratio_antitone}、
\leanref{TailResult}{tailBound_antitone} 中证明；
全域结论为 \leanref{TailResult}{tail_bound_all}。

\subsection{尾部素数定理}

\begin{theorem}[无限尾部的严格短区间]\label{tail:prime-theorem}
假设 H1、H2、H3、H6。对每个实数 $x>0$，若 $\log x\geq700$，
则存在自然数素数 $p$ 满足
\[
 x<p<x+4x^{3/4}.
\]
\end{theorem}

\begin{proof}
先取 $x=\e^L$，$L\geq700$，并使用式~\eqref{tail:parameters}。
将式~\eqref{tail:chjii} 的复数残差界与式~\eqref{tail:zero-bound} 相加，
三角不等式给出
\[
 \psi(y)-\psi(x)
 \geq h\bigl(1-Z_0-Z_1-Z_2-R\bigr).
\]
这里使用的是整个复数模，并未把可能的虚部误丢弃。
再由 $\theta(x)\leq\psi(x)$ 及引理~\ref{tail:arithmetic-errors}，
\[
 \theta(y)-\theta(x)
 \geq h\bigl(1-Z_0-Z_1-Z_2-R-P\bigr).
\]
式~\eqref{tail:global-budget} 说明括号内严格大于 $E$，从而
\[
 \theta(y)-\theta(x)>hE\geq\log y.
\]
应用式~\eqref{tail:open-extraction} 的逆否命题，得到严格开区间内的素数。
最后，对任意满足条件的 $x>0$ 令 $L=\log x$；
$\exp(\log x)=x$ 且 $h=4x^{3/4}$，结论随即成立。
\end{proof}

最终连接对应 \leanref{TailFromInputs}{tail_short_interval_exp} 和
\leanref{TailFromInputs}{tail_short_interval_from_inputs}。
后一个定理的外部参数恰是 H1、H2、H3、H6；
零点预算、单调性、素数幂界和严格端点提取均已由本项目或 Mathlib 的证明项给出。

\begin{remark}\label{tail:source-boundary}
对 CHJII 的源文核查确认了当前输入与 arXiv v3 陈述一致，相关定理、表格
和证明段落在 v2 与 v3 间未改；首行 $2.894$ 的构成及上游 $K$ 已作数值复算和
Arb 标量包围。正式期刊正文的精确版本对照及完整上游解析证明的独立重证仍未完成。
这些来源状态不影响以上条件推导的逻辑形式，也不能被上述 Lean 编译结果反向
解释为 H6 已在本库中获得证明。
\end{remark}

\section{整数下段：素数证书、CHL 十二行与指数拼接}
\label{sec:lower}

解析方法覆盖的是以实变量 $x$ 表示的巨大区间；要得到所有正整数 $n$
的结论，还必须填满 $x=n^6$ 的较小范围。本节在唯一外部输入 H7 下证明
\begin{equation}\label{lower:target}
 1\leq n\leq N_*:=1242317999994
 \quad\Longrightarrow\quad
 \exists p\text{ 为素数},\quad n^6<p<n^6+6n^5.
\end{equation}
其中 $1\leq n\leq1259$ 完全由本库的实际素性证明处理；
其余范围使用 H7 中明确列出的十二条 CHL 短区间结论。
本节不使用 QRH、CS、Rosser 或有限 RH。

这一区分对理解证明的可信基础很重要。历史数值文件曾记录一种有限基素性测试；
当前 Lean 证明的这些素数已另由递归 Lucas 证书证明，因此该测试及其适用范围
并不是当前定理的额外文献输入。另一方面，CHL 行仍是 H7，
不能因为它们的整数转换已经形式化，就声称 CHL 的上游解析定理也已形式化。

\subsection{1259 个小区间的确定素数见证}

对每个整数 $1\leq n\leq1259$，保存一个具体自然数 $p_n$。
形式化需要分别证明三件事：
\begin{equation}\label{lower:witness-obligations}
 \operatorname{Prime}(p_n),\qquad n^6<p_n,\qquad
 p_n<n^6+6n^5.
\end{equation}
后两项是整数算术命题，不能由“$p_n$ 大致在 $n^6$ 附近”代替。
前几项及末项如下，完整列表由随附数据与源文件提供：
\begin{center}
\begin{tabular}{r|r|r}
$n$ & $n^6$ & $p_n$\\\hline
1 & 1 & 2\\
2 & 64 & 67\\
3 & 729 & 733\\
4 & 4096 & 4099\\
5 & 15625 & 15629\\
1259 & 3982487126873086441 & 3982487126873086493
\end{tabular}
\end{center}
完整见证列表在 \texttt{results/lower\_range\_k6\_certificate.json} 中；
与其对应的定理是 \texttt{PrimeWitnesses.lean} 中的
\texttt{small\_prime\_1} 至 \texttt{small\_prime\_1259}。
每条定理引用一个已经证明的素性命题，并直接证明两个严格整数不等式。

\begin{lemma}[本项目使用的 Lucas 素性准则]\label{lower:lucas}
设整数 $P\geq2$，并已知 $P-1$ 的完整素因数分解。
若存在模 $P$ 的元素 $a$，满足
\begin{equation}\label{lower:lucas-congruences}
 a^{P-1}\equiv1\pmod P,
 \qquad
 a^{(P-1)/q}\not\equiv1\pmod P
 \quad\text{对每个素因子 }q\mid P-1,
\end{equation}
则 $P$ 为素数。
\end{lemma}

\begin{proof}
第一个同余保证 $a$ 是模 $P$ 的可逆元，其阶整除 $P-1$。
如果该阶是 $P-1$ 的真因子，则 $(P-1)/\operatorname{ord}(a)$
有某个素因子 $q$，从而该阶整除 $(P-1)/q$，与第二个同余条件矛盾。
因此 $a$ 的阶恰为 $P-1$。由有限单位群的阶关系，
$P-1\leq\varphi(P)$；而 $P\geq2$ 时 $\varphi(P)\leq P-1$，
等号只可能在 $P$ 为素数时成立。
\end{proof}

实际实现使用 Mathlib 的 \texttt{lucas\_primality}。
给出 $P-1=\prod_iq_i^{e_i}$ 时，既检查乘法恒等式，
也递归引用每个 $q_i$ 的已证明素性；然后对所有这些素因子检查
式~\eqref{lower:lucas-congruences}。由于 $q_i<P$，依赖关系严格向较小整数
下降，形成可复用节点的有向无环证书。完整数据共有
\emph{5349 个素数节点}，分为 27 个源文件
\texttt{Pratt/Block000.lean} 至 \texttt{Pratt/Block026.lean}。
小于 $1000$ 的节点由 \texttt{norm\_num} 的有限算术素性证明处理，
其余节点用上述 Lucas 准则。

例如，对于见证 $P=15629$，采用
\[
 P-1=2^2\cdot3907,\qquad a=2.
\]
在已经证明 $2$ 与 $3907$ 为素数之后，需检查
\[
 2^{15628}\equiv1,\qquad
 2^{7814}\equiv15628\ne1,\qquad
 2^4\equiv16\ne1\pmod{15629}.
\]
这条证书可在 \leanref{Pratt/Block005}{prime_15629} 中直接核查。
其它较大的见证遵循同一逻辑，不要求读者相信一个外部程序输出的
\texttt{is\_prime=true}。

Python/FLINT 负责寻找因数分解和模幂见证并生成候选 Lean 源码。
这些候选只有在整数乘法、模幂和递归素性证明均通过内核检查后才能使用。
模幂计算由 \texttt{reduce\_mod\_char} 生成证明项；这里没有采用
\texttt{native\_decide} 作为捷径，也没有把 Python 的返回值注册成数学公理。

\begin{proposition}[小整数范围]\label{lower:small}
对每个 $1\leq n\leq1259$，存在素数 $p$ 满足
$n^6<p<n^6+6n^5$。
\end{proposition}

\begin{proof}
对有限个整数 $n$ 分情形，每个叶节点使用已经证明的
式~\eqref{lower:witness-obligations}。
实际源文件用二分情形树组织这些叶节点，以控制证明深度；
二分组织方式不改变被穷尽的整数集合。
\end{proof}

对应定理为 \leanref{PrimeWitnesses}{small_range_prime}。
特别地，$n=1$ 由 $p=2$ 处理，故后面要求 $n\geq2$ 的合数端点论证
不会遗漏这个最小正整数。

\subsection{CHL 的区间方向与严格端点转换}

H7 的每一行采用下列准确的方向和端点约定：给定 $\Delta>1$ 与 $X$，
\begin{equation}\label{lower:CHL-input}
 y\geq X\quad\Longrightarrow\quad
 \exists p\text{ 为素数},\quad
 y\left(1-\frac1\Delta\right)<p\leq y.
\end{equation}
源文为 \cite{CHL,CHLcorr} 的修订结论及指定表行。
这是左开右闭的区间；不能直接将右端点的不严格不等式抄成严格不等式。
接口 \leanref{Definitions}{CHLInterval} 保留了这一差别。

\begin{lemma}[单行 CHL 到六次幂短区间]\label{lower:CHL-conversion}
假设式~\eqref{lower:CHL-input} 成立。若整数 $n\geq2$ 满足
\begin{equation}\label{lower:row-conditions}
 X\leq n^6,\qquad n\leq6(\Delta-1),
\end{equation}
则存在素数 $p$，使 $n^6<p<n^6+6n^5$。
\end{lemma}

\begin{proof}
令 $y=n^6+6n^5=n^5(n+6)$。
因 $y\geq n^6\geq X$，可对 $y$ 使用 H7 的对应行。
关键的左端点比较为
\begin{align}
 y\left(1-\frac1\Delta\right)-n^6
 &=\frac{n^5}{\Delta}\bigl(6(\Delta-1)-n\bigr)\geq0.
 \label{lower:left-comparison}
\end{align}
于是源文给出的素数满足
$p>y(1-1/\Delta)\geq n^6$。
右端仅有 $p\leq y$，但 $n\geq2$ 时两个整数因子
$n^5$ 和 $n+6$ 都严格大于 $1$，所以 $y=n^5(n+6)$ 是合数。
由于 $p$ 是素数，$p\ne y$，因此 $p<y$。
\end{proof}

这两个步骤分别对应
\leanref{CHL}{nat_sixth_short_of_CHL} 与
\leanref{Reduction}{sixth_right_endpoint_not_prime}。
注意式~\eqref{lower:row-conditions} 容许等号：即使
$n=6(\Delta-1)$，左端仍由 CHL 的严格不等式保证 $p>n^6$。

\subsection{全部十二行及整数阈值}

下表列出当前 \leanref{CHLRange}{CHLRows} 的全部输入行，
并给出本项目使用的整数覆盖区间 $[a_j,b_j]$。
这里 $a_j$ 是已核验的安全下阈值，不声称是最小可能阈值；
$b_j=6(\Delta_j-1)$ 是精确整数。

\begin{center}
\small
\begin{tabular}{c|r|r|r}
$X_j$ & $\Delta_j$ & $a_j$ & $b_j=6(\Delta_j-1)$\\\hline
$4\cdot10^{18}$ & $39097000$ & $1260$ & $234581994$\\
$\e^{43}$ & $41816800$ & $1296$ & $250900794$\\
$\e^{46}$ & $163940000$ & $2136$ & $983639994$\\
$\e^{50}$ & $1061200000$ & $4161$ & $6367199994$\\
$\e^{55}$ & $10288400000$ & $9573$ & $61730399994$\\
$\e^{60}$ & $76918400000$ & $22027$ & $461510399994$\\
$\e^{75}$ & $174043000000$ & $268338$ & $1044257999994$\\
$\e^{90}$ & $184304000000$ & $3269018$ & $1105823999994$\\
$\e^{105}$ & $191886000000$ & $39824785$ & $1151315999994$\\
$\e^{120}$ & $197917000000$ & $485165196$ & $1187501999994$\\
$\e^{135}$ & $202553000000$ & $5910522064$ & $1215317999994$\\
$\e^{150}$ & $207053000000$ & $72004899338$ & $1242317999994$
\end{tabular}
\end{center}

每行使用两类不同的证据。$\Delta_j$ 和源域 $X_j$ 是 H7 的文献输入；
而将其转换为 $[a_j,b_j]$ 所需的
\begin{equation}\label{lower:threshold-checks}
 X_j\leq a_j^6,\qquad b_j\leq6(\Delta_j-1)
\end{equation}
均在本库中证明，未再假设一个整数区间结论。
第一行的下界是直接的整数计算：
\[
 1260^6=4001504141376000000>4\cdot10^{18}.
\]
其余十一条指数不等式完整列为
\begin{alignat}{2}
 \e^{43}&\leq1296^6,&\qquad \e^{46}&\leq2136^6,\notag\\
 \e^{50}&\leq4161^6,& \e^{55}&\leq9573^6,\notag\\
 \e^{60}&\leq22027^6,& \e^{75}&\leq268338^6,\notag\\
 \e^{90}&\leq3269018^6,& \e^{105}&\leq39824785^6,\notag\\
 \e^{120}&\leq485165196^6,&
 \e^{135}&\leq5910522064^6,\notag\\
 \e^{150}&\leq72004899338^6.&&
 \label{lower:eleven-exp-thresholds}
\end{alignat}
这些比较由 \texttt{Thresholds.lean} 中
\texttt{chl\_threshold\_43} 至 \texttt{chl\_threshold\_150}
的反射式二进区间证明验证，设置的工作精度为 40 位。
指数的数值包围本身生成内核可检查的证明，
并非先在外部取 $\lceil\e^{j/6}\rceil$ 再把这个整数当作正确答案。

如果 $a_j\leq n\leq b_j$，幂函数在非负实轴的单调性给出
$X_j\leq a_j^6\leq n^6$，同时 $n\leq6(\Delta_j-1)$。
因此引理~\ref{lower:CHL-conversion} 对该行覆盖的每个整数都适用。
形式化把这个共同步骤封装成
\leanref{CHLRange}{nat_short_of_CHL_integer_row}。

\begin{proposition}[CHL 行之间没有整数空隙]\label{lower:CHL-coverage}
在 H7 下，对每个整数
$1260\leq n\leq1242317999994$，存在素数 $p$ 满足
$n^6<p<n^6+6n^5$。
\end{proposition}

\begin{proof}
按表中顺序记十二个区间为 $[a_j,b_j]$，$0\leq j\leq11$。
表内精确整数满足
\[
 a_0=1260,\qquad b_0<b_1<\cdots<b_{11}=N_*,
 \qquad a_{j+1}\leq b_j+1\quad(0\leq j<11).
\]
所以 $[a_0,b_0]$ 与其后逐行增加的区间覆盖整个整数区间
$[1260,N_*]$。等价地，对给定的 $n$，取第一个满足 $n\leq b_j$ 的行；
若 $j>0$，前一行未覆盖便给出 $n\geq b_{j-1}+1\geq a_j$，
从而这一行的上下条件都成立。
每个比较都可从表中整数直接验证；无需关于素数分布的进一步输入。
\end{proof}

对应定理为 \leanref{CHLRange}{chl_integer_coverage}。
与命题~\ref{lower:small} 拼接即得
\leanref{LowerRange}{lower_range_complete}，也就是式~\eqref{lower:target}。
这里的 $1259/1260$ 接口是连续的；我们没有以“有限多个剩余例外可计算”
结束论证，而是给出了这些例外的实际素数证明。

\subsection{与 \texorpdfstring{$L=167$}{L=167} 的精确拼接}

要与后面的实变量中段相接，还需证明式~\eqref{lower:target} 至少覆盖
所有满足 $\log(n^6)\leq167$ 的正整数。这依赖一条另外验证的超越不等式：
\begin{equation}\label{lower:join-exp}
 \e^{167}<(1242317999994)^6=N_*^6.
\end{equation}
\leanref{Thresholds}{lower_join_exp} 使用 20 位反射式二进区间证明建立该式。
它与上一组十一条阈值检查一样，是 Lean 中的实数不等式证明。

\begin{lemma}[指数接口]\label{lower:log-join}
设 $n\in\N$ 且 $n\geq1$。若 $\log((n:\R)^6)\leq167$，则 $n\leq N_*$。
\end{lemma}

\begin{proof}
由 $n>0$，指数函数的单调性给出
\[
 n^6=\exp(\log(n^6))\leq\e^{167}<N_*^6.
\]
若 $n>N_*$，非负整数的六次幂保持序关系，将与最后一个严格不等式矛盾。
所以所需的非严格上界 $n\leq N_*$ 成立。
\end{proof}

\begin{theorem}[完整整数下段]\label{lower:log-theorem}
在 H7 下，对每个 $n\in\N$，若 $n\geq1$ 且
$\log((n:\R)^6)\leq167$，则存在素数 $p$ 满足
\[
 n^6<p<n^6+6n^5.
\]
\end{theorem}

\begin{proof}
先由引理~\ref{lower:log-join} 得 $n\leq N_*$，再使用
小范围证书与 CHL 行已经拼成的式~\eqref{lower:target}。
\end{proof}

对应实现为 \leanref{CHLRange}{lower_join_of_log} 和
\leanref{LowerRange}{lower_log_range}。
式~\eqref{lower:join-exp} 为严格不等式，说明两个范围有重叠；
并不需要猜测 $\e^{167/6}$ 位于哪两个相邻整数之间。

\subsection{H7 的来源状态与本节证明的边界}

CHL 原文与勘误均有正式期刊发表记录。当前十二行与作者 arXiv v3 修订本的
数值、定义域以及左右端点约定相符；但是，这与独立验证其全部上游计算不是同一件事。
源文复核发现，公开的有限位辅助参数在
$\log X=50,60,75,90,105,135$ 六行生成的 $\Delta$
略小于表中目标值。该方向差异已经用以印刷十进制数为有理输入的 Arb 计算确认。
由于较大的 $\Delta$ 对应较短的区间，不能用略小的生成值直接证明表中较强的行。
小幅调整辅助参数后能够达到目标的数值观察，也不代替完整严格认证。

本节因此明确接受的是 H7 中的十二条\emph{区间命题}，
不是声称已由公开舍入见证在本库重证这些命题。
本节对它们之后的步骤——整数阈值、各行覆盖、素数证书、严格端点和指数拼接——
给出了完整证明。正式勘误正文的精确表格比对以及 CHL 判据的完整独立认证仍列在来源审查中。
这一限制应与定理同时阅读，而不能被“所有下游模块编译通过”一句话覆盖。

\section{全范围拼接与六次幂结论}
\label{sec:assembly}

本节将各段的实际素数定理拼成关于所有正整数的结论。
需要区分两个任务：前面各节已经证明具体范围中的素数存在性；本节再证明
这些范围覆盖全部正整数，并核对实数幂、整数幂及所有严格端点。
这一步没有再引入零点估计或数值假设。

为与形式化接口一致，记
\begin{align}
 \mathcal S_6(x)&:\quad
   \exists p\in\N,\quad p\text{ 为素数},\quad x<p<x+6x^{5/6},
   \label{assembly:S6}\\
 \mathcal S_4(x)&:\quad
   \exists p\in\N,\quad p\text{ 为素数},\quad x<p<x+4x^{3/4},
   \label{assembly:S4}\\
 \mathcal N_6(n)&:\quad
   \exists p\in\N,\quad p\text{ 为素数},\quad n^6<p<n^6+6n^5.
   \label{assembly:N6}
\end{align}
前两个定义中的比较在 $\R$ 中进行，自然数 $p$ 通过标准嵌入视为实数；
第三个定义中的幂与不等式都是自然数运算。
对应定义分别是 \leanref{Definitions}{ShortInterval6}、
\leanref{Definitions}{ShortInterval4} 和
\leanref{Definitions}{NatSixthShort}。
下标 $4$ 和 $6$ 在这里表示选用的短区间长度；
它们本身不是某个全范围幂次猜想已经成立的声明。

\subsection{已经获得的各段及其精确依赖}

所有实变量段都以 $x>0$、$L=\log x$ 为前提。
下表把用于最终拼接的五个结论放在一起。
其中的每个区间端点均包含在相应定理中。

\begin{center}
\small
\begin{tabular}{c|c|c|l}
范围 & 变量及条件 & 结论 & 外部输入\\\hline
整数下段 & $n\geq1$, $\log(n^6)\leq167$ & $\mathcal N_6(n)$ & H7\\
下中段 & $167\leq L\leq715/4$ & $\mathcal S_6(x)$ & H1, H2, H3, H4, H5\\
上中段 & $715/4\leq L\leq400$ & $\mathcal S_6(x)$ & H1, H2, H3, H4\\
有限远段 & $400\leq L\leq700$ & $\mathcal S_4(x)$ & H1, H2, H3, H4\\
无限尾部 & $700\leq L$ & $\mathcal S_4(x)$ & H1, H2, H3, H6
\end{tabular}
\end{center}

整数下段是第~\ref{sec:lower} 节的定理~\ref{lower:log-theorem}。
下中段与上中段分别由
\leanref{M3LowerResult}{lower_middle_short_interval_from_inputs} 和
\leanref{M3MiddleUpper}{middle_upper_short_interval_from_inputs} 给出。
有限远段为
\leanref{M3FarUpper}{far_short_interval_from_inputs}，
无限尾部为第~\ref{sec:tail} 节的定理~\ref{tail:prime-theorem}。
中段分点 $715/4=178.75$ 只决定采用哪一个已证明的解析估计，
不会在最终定理中留下这个分点处的额外假设。

以上依赖表也说明，H5 的使用仅在下中段，H6 仅在无限尾部，H7 仅在整数下段。
最终定理统一接受七个输入，是因为它同时覆盖这些范围；
这不表示每一个局部定理都使用了全部七个输入。

\subsection{四次长度与六次长度的比较}

\begin{lemma}[短区间长度的包含关系]\label{assembly:length-comparison}
若 $x\geq1$，则
\begin{equation}\label{assembly:length-inequality}
 4x^{3/4}\leq6x^{5/6}.
\end{equation}
因此 $\mathcal S_4(x)$ 蕴含 $\mathcal S_6(x)$。
\end{lemma}

\begin{proof}
因为 $3/4\leq5/6$ 且 $x\geq1$，实数幂关于指数的单调性给出
$x^{3/4}\leq x^{5/6}$。后者非负，故
\[
 4x^{3/4}\leq4x^{5/6}\leq6x^{5/6}.
\]
若 $p$ 是 $\mathcal S_4(x)$ 的素数见证，则
\[
 x<p<x+4x^{3/4}\leq x+6x^{5/6}.
\]
严格不等式与非严格不等式的传递仍给出严格的最终右界。
\end{proof}

相应实现为 \leanref{Reduction}{four_length_le_six_length} 和
\leanref{Reduction}{short_six_of_short_four}。
这里必须说明 $x\geq1$；当 $0<x<1$ 时，单凭指数大小不能作同方向比较。
在本证明的有限远段和无限尾部，$L\geq400$，所以这个条件自动满足；
在最终应用 $x=n^6$ 中，它也由 $n\geq1$ 直接得到。

\begin{proposition}[统一的解析上段]\label{assembly:analytic-range}
在 H1 至 H6 下，对每个 $x>0$，若 $\log x\geq167$，则
$\mathcal S_6(x)$ 成立。
\end{proposition}

\begin{proof}
记 $L=\log x$。若 $L\leq400$，则在
$L\leq715/4$ 与 $L>715/4$ 两种情况下，分别使用下中段与上中段的
$\mathcal S_6(x)$ 结论。
若 $L>400$，则 $x>1$；当 $L\leq700$ 时使用有限远段，
当 $L>700$ 时使用无限尾部。后两种情况先得到 $\mathcal S_4(x)$，
再由引理~\ref{assembly:length-comparison} 得到 $\mathcal S_6(x)$。
每种情形均由某个闭区间定理覆盖，故 $167$、$715/4$、$400$、$700$
这四个接口没有遗漏。
\end{proof}

中段的具体合并见
\leanref{FinalResult}{middle_short_interval_from_inputs}。
把三个已构造的解析段交给
\leanref{Assembly}{real_short_above_167_of_segment_inputs}
便得到命题~\ref{assembly:analytic-range} 的形式化版本。

\subsection{从实变量区间回到整数六次幂}

\begin{lemma}[实数幂与自然数幂的精确转换]\label{assembly:power-conversion}
对每个自然数 $n$，有
\begin{equation}\label{assembly:rpow-identity}
 \bigl((n:\R)^6\bigr)^{5/6}=(n:\R)^5.
\end{equation}
若 $\mathcal S_6((n:\R)^6)$ 成立，则 $\mathcal N_6(n)$ 成立。
\end{lemma}

\begin{proof}
自然数在实数中的像非负，故非负实底数的幂法则适用。
由 $6(5/6)=5$，得到式~\eqref{assembly:rpow-identity}。
将这个等式代入 $\mathcal S_6$ 的右端点后，见证素数满足
\[
 (n:\R)^6<(p:\R)<(n:\R)^6+6(n:\R)^5.
\]
自然数到实数的嵌入保持加法、乘法、整数幂和严格序关系，
因此两条实数不等式恰给出式~\eqref{assembly:N6} 的两条自然数不等式。
素数 $p$ 本身不发生改变。
\end{proof}

实现为 \leanref{Reduction}{sixth_rpow_five_sixths} 与
\leanref{Reduction}{nat_short_of_real_short}。
这一转换没有使用四舍五入，也不需要对实数端点先取整。

\begin{lemma}[相邻六次幂之间的严格余量]\label{assembly:binomial-gap}
对每个自然数 $n$，
\begin{equation}\label{assembly:binomial}
 n^6+6n^5<(n+1)^6.
\end{equation}
\end{lemma}

\begin{proof}
二项式展开给出
\[
 (n+1)^6-(n^6+6n^5)
 =15n^4+20n^3+15n^2+6n+1>0.
\]
各幂非负，而常数项为 $1$，所以严格正性甚至对 $n=0$ 也成立。
\end{proof}

对应定理是 \leanref{Reduction}{sixth_binomial_gap}。
不过最终的素数存在性定理仍只声明 $n\geq1$：$n=0$ 时区间 $(0,1)$
没有素数，不能因为这个单独的二项式比较对 $n=0$ 成立便扩大主定理范围。

\subsection{全范围条件定理}

\begin{theorem}[七个明确输入下的六次幂短区间定理]
\label{assembly:main-short}
假设 H1 至 H7 成立。对每个自然数 $n\geq1$，存在自然数素数 $p$，使
\begin{equation}\label{assembly:strong-chain}
 n^6<p<n^6+6n^5<(n+1)^6.
\end{equation}
\end{theorem}

\begin{proof}
固定 $n\geq1$，令 $x=(n:\R)^6>0$，$L=\log x$。
若 $L\leq167$，由定理~\ref{lower:log-theorem} 直接得到
$\mathcal N_6(n)$。
若 $L>167$，则命题~\ref{assembly:analytic-range} 给出
$\mathcal S_6(x)$，再由引理~\ref{assembly:power-conversion} 得
$\mathcal N_6(n)$。
两种情形已经穷尽所有 $n\geq1$，且各自均给出实际的自然数素数见证。
最后附上引理~\ref{assembly:binomial-gap} 的严格比较，得到完整链式不等式。
\end{proof}

\begin{corollary}[六次幂之间的素数]\label{assembly:sixth-powers}
在 H1 至 H7 下，对每个自然数 $n\geq1$，开区间
$(n^6,(n+1)^6)$ 中至少有一个素数。
\end{corollary}

这两个结论分别由
\leanref{FinalResult}{sixth_power_short_from_named_inputs} 及
\leanref{FinalResult}{sixth_powers_from_named_inputs} 实现。
后一个定理的目标
\leanref{Definitions}{SixthPowerConclusion}
直接量化所有 $n\in\N$ 并要求 $1\leq n$；它不是“所有已检查整数”的缩写。
前一个定理保留了更强的内端点 $n^6+6n^5$，所以同时记录了
式~\eqref{assembly:strong-chain} 中的额外余量。

\subsection{中间接口已经构造，而非留下局部结论作为前提}

证明库保留两个便于复用的中间接口。
\leanref{Assembly}{AnalyticSegmentInputs} 的三个字段分别陈述
中段、有限远段和尾部的素数结论；
\leanref{AssemblyWithTail}{FiniteAnalyticSegmentInputs}
只包含前两个有限段。
如果仅证明“给定这些字段，就能推出六次幂结论”，那仍然只是一个拼接定理，
因为困难的局部素数结论仍作为未证明字段存在。

当前最终定理补齐了这一步。
\leanref{FinalResult}{finite_analytic_segments_from_inputs}
把已证明的中段与有限远段定理填入两个字段；
\leanref{TailFromInputs}{tail_short_interval_from_inputs}
再提供尾部字段。
\leanref{FinalResult}{sixth_power_short_from_named_inputs}
和 \leanref{FinalResult}{sixth_powers_from_named_inputs}
的参数表因此只留下 H1 至 H7，并不再要求调用者提供
一个局部素数结论、一个单元最大值不等式或一个未证明的解析段接口。

这一点还可以在历史单元的对应关系中直接检查。
\leanref{FinalCellCoverage}{all_94_cells_actual_prime_intervals}
证明每个保存的中段细分单元内都有 $\mathcal S_6(x)$；
\leanref{FinalCellCoverage}{all_60_far_cells_actual_prime_intervals}
证明每个保存的远段单元内都有 $\mathcal S_4(x)$。
它们先在内核中检查每个单元落入相应的全段，再调用实际的素数定理。
因此保存的有理数最大值并不是这两个结论的额外解析前提。
同时，当前远段的形式化使用梯形权函数。
它覆盖历史六十个单元的目标区间，但不等于把历史多项式权函数手稿逐行翻译为 Lean。
这两种研究记录应分别保留，避免将“目标范围已形式化”误写成
“每一种历史证明实现均已形式化”。

\begin{remark}[主定理的逻辑范围]\label{assembly:scope}
本文完成的是七个明确数学输入之后的全范围推导。
H1 是本路线承认的 QRH 假设；H2 至 H7 是精确定义的文献命题接口，
其来源状态及独立审查程度须分别陈述。
Lean 证明这些命题一旦成立，式~\eqref{assembly:strong-chain}
便对每个 $n\geq1$ 成立；它不凭编译成功而证明七个输入本身。
特别地，这不是只以 QRH 一项为参数的库内定理，也不是无条件六次幂定理。
至于 $\mathcal S_4(x)$，本文已经在 $\log x\geq400$ 的实变量范围证明它，
但没有据此宣称所有正整数的四次幂猜想已经解决。
\end{remark}

\section{方法的可复用部分与下一步问题}\label{sec:methods}

\subsection{从定位损失到选择方法}
研究的起点是把一个全范围问题拆成可测量的若干损失：零点贡献、截断误差、非素数素数幂以及端点贡献。若归一化后这些量的严格上界之和小于 $1$，主项就留下正的素数贡献。这个框架的作用在于辨认困难来自哪里；提高小数精度只能减少数值包围误差，不能缩小一个本来很大的解析上界。

早期围绕 $k=4$ 的试验揭示了中段与尾段的区别。尾段可以利用随 $L$ 衰减的指数因子建立解析单调性；中段则对密度常数、权函数的变换衰减和有限高度的拆分更加敏感。因而本项目把 $k=6$ 作为能够闭合全部范围的阶段成果，同时保留对较小指数的局部计算和结构分析。

梯形权的选择兼顾两个要求：权非负且在严格开区间之外为零，使正加权和能提取素数；其四点差分表达又允许从同一个积分显式公式得到零次、一次和二次三种高度衰减界。取三者的最小值，可以在不同高度利用合适的估计。更高平滑阶数并不自动更好，因为主项质量、导数范数和相对区间长度必须同时计算。

\subsection{形式化促成的数学重组}
原始书面证明逐个认证中段和远段单元；远段还采用另一种多项式权。最终形式化把主要有限范围重组为两个基准证书、一个连续传播定理和一个有限高度 RH 比较。这样既减少重复的超越估计，也给出了对所有实参数成立的结构性解释。

连续传播的关键不是把一个成功参数复制到邻近采样点，而是同时平移 $A$ 和 $u$、缩放 $q$。核的高度质量保持不变，密度包围则由导数上界控制为不增加。第\ref{sec:coverage}节证明的适用条件包含 $k\ge4$；这说明该引理本身能服务于多个指数，但全范围素数结论还要求某个合法起点具有严格小于 $1$ 的完整预算，并与整数下段衔接。

有限高度 RH 的使用也有明确目的。在 $L$ 向下移动时，一般零点密度预算会变坏；低于 $H=3\cdot10^{12}$ 的零点却可以使用 $\beta=1/2$ 的精确信息。将低高度单独扣出后，高高度核的缩放补偿了部分损失，留下可由三档有理系数认证的比较。

\subsection{计算余量应如何解释}
预算小于 $1$ 的差额是所用估计的确定性余量。例如下中段的保守上界为 $0.966283$，因而这一预算还留有 $0.033717$。这表示在当前分项和证明条件下，扣除项严格不足主项；它不是命题正确的概率，也不允许把每个上游参数任意改动同一百分比。

原始 $94$ 单元证书另验证过全部三列 CS 系数同时增加 $6\%$ 的敏感性试验。该结果属于原始证书与其固定包围。新下中段采用另一套保守比较，故应分别报告其实际预算，不能把原敏感性结论直接迁移给当前 Lean 定理。研究历史保留两套记录及其明确对应的数学版本。

CHL 的舍入参数复现问题同样要区分目标常数与见证参数。我们在 H7 中接受的是指定短区间结论；作者表中有限位的辅助参数是建立它的见证。六行印刷参数生成的 $\Delta$ 略低于报告目标，微调参数后的普通数值试验显示仍有可行余量。严格区间诊断已经排除了浮点伪象，但完整判据的独立认证仍是一项清晰、有限的后续工作。这项问题保留在输入登记中，而不改变本次给定 H1--H7 的逻辑推导。

\subsection{通向较小指数的任务}
若考虑一般长度 $h=kx^{1-1/k}$，则控制核的参数具有 $q\asymp x^{1/k}/k$ 的尺度。减小 $k$ 会增大相对所需截断高度，并使同一 $L$ 下的加权零点预算更难压低。可复用的部分包括权函数恒等式、密度分层、高度尾项、整数矩认证以及部分连续传播引理；需要重新求解的部分包括基准参数、全区间预算、下端覆盖与两者之间的接合。

因此，对 $k=4$，下一步应寻找能够与已知整数下段相接的有效基准区间，并定位阻碍预算下降的密度区间与高度区间。对 $k=3$，现有连续传播定理的 $k\ge4$ 条件本身还需改进。局部成功和一般形式的引理都是可积累成果，只有所有范围闭合后才升级为新的全范围结论。

本文的已证明成果、历史探索及未来任务在仓库中分层保存。这样的记录允许后续研究复用现有组件，而不把某次未闭合的探索误读为主定理的一部分。

\section{形式化对应、复现与研究贡献}\label{sec:formalization}

\subsection{最终形式化命题}
本项目使用 Lean 的实际自然数素性谓词与 Mathlib 的实际黎曼 zeta 函数。最终定理为
\begin{Verbatim}[fontsize=\small,breaklines=true]
theorem sixth_powers_from_named_inputs
  (HQ : QRHSevenEighths) (HC : CSv2)
  (HR : RosserZeroCount) (HI : IntegratedExplicitFormula)
  (HFR : FiniteRH3e12) (HJ : CHJIIShortInterval)
  (HCHL : CHLRows) : SixthPowerConclusion
\end{Verbatim}
其中 \repofile{SixthPowerConclusion} 展开即为对所有自然数 $n\ge1$，存在自然数素数 $p$ 严格位于相邻六次幂之间。更强的 \repofile{sixth_power_short_from_named_inputs} 给出 $p<n^6+6n^5$。

七项输入在定理参数中显式出现。其后没有再假定“中段成立”“某个数值证书正确”或“加权零点和满足所需上界”。这些中间结论由前述数学构造及项目中的证明项给出。每章的形式化对应，以及仓库的证明导读，将人类表述连接到实际声明。声明名存在只能确认链接有效，数学含义的一致性仍须通过陈述与证明的逐项核对。

\subsection{证书生成与证书检查}
本项目的有限计算有三种不同角色。
\begin{enumerate}
\item 素性见证给出因数分解、模幂和递归 Lucas/Pratt 数据，Lean 检查这些数据确实证明每个所需整数为素数。
\item 解析证书提供有理包围与有限整数矩。Lean 证明真实指数、对数及实幂落在这些包围内，再检查整数不等式。Python 或 Arb 的候选输出可以被更换，只要重新提交的数据仍通过同一检查。
\item 文献来源审查中的 Arb\cite{Arb} 重放核验指定数值公式，其结果作为外部证据保存。它不自动成为 Lean 定理，也不能替代尚未重证的上游解析命题。
\end{enumerate}
这一划分解释了为什么仓库既保留数据，也保留生成器和证明理论：生成器使数据可重现，证明项说明数据为什么足够，实际内核检查确认指定数据满足形式化命题。

\subsection{固定环境与验收记录}
当前源码固定 Lean \repofile{v4.36.0-rc2}\cite{LeanPinned}，Mathlib\cite{MathlibPinned} 提交为
\begin{center}\small\nolinkurl{9d581dfff2e41de989072ea63abf3daf69311520}\end{center}
全部传递依赖由 \repofile{lean/lake-manifest.json} 记录。历史最终验收在一个原先不存在的输出目录中重建了 $241$ 个项目模块，全部通过，并在对应产物上执行最终依赖审计。该记录位于 \repofile{audit/builds/final-20261106/}，其源码与产物散列由验收检查点绑定。

依赖审计同时打印最终定理的完整参数和证明项中实际使用的本地声明。最终公理报告仅含 \repofile{propext}、\repofile{Classical.choice}、\repofile{Quot.sound}。由于外部命题作为参数出现，单独观察公理列表不能说明它们已经得到证明；完整类型与依赖记录必须一起阅读。

本次整理保留上述数学源码与原验收记录，同时在新的干净源码导出上执行独立输出目录的重建和原始证书重放。完成状态、执行环境、散列与耗时记录于 \repofile{audit/archive_review_2026-11-06.json}。它是整理版本的复现证据，不改写旧验收发生时的记录。

\subsection{读者的复现顺序}
从指定发行提交检出仓库后，先检查发布文件散列；再按 \repofile{REPRODUCING.md} 安装固定环境，验证原始证书和生成源码；随后重建全部 Lean 模块并运行最终依赖审计。论文在 \repofile{paper/current/} 中单独构建。排版编译只产生阅读版本，数学验证由上述证明和数值检查完成。

发布包包含本项目全部 Lean 源码，包括生成的素性与数值证书证明文件。Mathlib 等第三方依赖按固定提交取得，缓存和机器专用产物不作为原创源码分发。原始证书、历史报告与本次整理后的文档之间通过散列和版本索引对应。

\subsection{AI 研究贡献与审阅范围}
按项目发起人的署名约定，Claude 与 Codex 作为 AI 系统贡献者列名。Codex 负责 QRH 路线的研究组织、解析推导、数值实现、证书细化、Lean 实现、来源核查及本稿整理。Claude 负责传统显式估计路线，并对 QRH 路线早期版本提供审阅；其意见推动了 CS v2 依赖显式化及中段网格细化。那份历史审阅的范围限于对应版本。

项目发起人提出目标、决定路线与质量要求、要求公开完整材料，并审核发布。人类姓名、单位与联系信息暂未填写。Codex 的不同代理上下文参与实现和交叉检查，相关报告区分实现与复核角色；这些记录不等同于人类专家同行评审。

部分交互、模型后端标识及随机性信息未完整保存，故研究过程按可核对的工作记录、源码和数学产物归档。可重复执行的是固定数学输入、程序、证书生成和指定工具链下的验证。完整贡献说明见 \repofile{AI_USE.md} 与根目录 \repofile{AUTHORS.md}。

QRH 路线原创代码采用 0BSD，原创论文、说明和数据采用 CC0。被引用论文、软件依赖与作者原始代码保留各自许可。参考文献列出正式出版信息和精确版本；来源下载及历史版本定位信息另见输入登记与专项报告。

\appendix
\section{证明版本与历史材料的对应}\label{app:versions}

本附录帮助读者区分当前主证明与仍有研究价值的历史材料。历史文件保留其原文及当时结论；当前状态由新主稿和验收索引给出。

\begingroup\small
\begin{longtable}{p{0.20\linewidth}p{0.34\linewidth}p{0.34\linewidth}}
\toprule
事项 & 历史版本 & 当前论文对应\\
\midrule\endfirsthead
\toprule 事项 & 历史版本 & 当前论文对应\\\midrule\endhead
远段权函数 & 原稿使用多项式权及 $60$ 个单元证书 & 第\ref{sec:trapezoid}--\ref{sec:coverage}节采用梯形权基准证书与连续传播\\
中段组织 & 精化后的 $94$ 个闭单元逐一认证 & 基准区间、连续传播和有限高度 RH 比较；另有 Lean 定理覆盖历史所有单元\\
数值余量 & 最坏旧单元预算约 $0.9397359968$；原证书支持统一 $6\%$ 系数试验 & 当前三类完整预算为 $0.939772$、$0.966283$、$0.450022$，各自按实际适用范围使用\\
初始形式化 & 有限见证、部分数值与装配引理；仍有局部解析参数 & 当前最终定理仅保留 H1--H7，所有项目局部解析与数值义务已经供给证明\\
积分公式来源 & Hildebrand 讲义作为直接出处 & 正式期刊论文接受稿的精确公式，详见 H4\\
CHJII 与 CHL 来源核查 & 采用固定作者稿 & 后续报告补充版本对照、常数重放及 CHL 舍入见证诊断，详见 H6、H7\\
\bottomrule
\end{longtable}\endgroup

对于复用历史结果的读者，关键是按其原来的参数、来源版本和检验范围解释记录。例如，旧多项式证书仍是原书面链的有效计算材料，却不是当前远段 Lean 证明使用的权函数；原 $6\%$ 试验不能仅因出现在同一仓库中，就归给新下中段比较。

\section{材料索引与持久标识}
当前阅读入口为 \repofile{README.zh-CN.md} 和 \repofile{documentation/README.md}；前提登记为 \repofile{documentation/inputs.md}；命题对应为 \repofile{documentation/blueprint/}；研究历史为 \repofile{research-history/}；完整复现为 \repofile{REPRODUCING.md}。原始论文位于上一层 \repofile{paper/main.pdf}，并在其索引中明确标为历史稿。

每次正式发布以 Git 提交及标签固定全部内容，\repofile{MANIFEST.json} 固定路线包文件散列，原验收检查点固定其对应的源码与构建证据。用户审核后的发布说明应给出实际提交标识，不在待审核稿中预填尚未发生的发布编号。

\clearpage
\bibliographystyle{plainurl}
\bibliography{references}
\end{document}
